\chapter{Matrix branes, relations and central rank}
\label{mht:ch:matrix-branes}

The relation resolving the commutative plane becomes a matrix
moment equation. Its scalar component vanishes identically, so
the relation complex retains a negative-degree class at every
rank. Ordinary invariant functions forget this class. Gauge
cochains add further states. These differences determine the
recipient of a Hamiltonian action.

We first compute the ordinary trace algebra and its stable
coordinate. We then quantize matrix entries and obtain a global
ordinary reduction with a strict Moyal trace map. The full BRST
complex requires additional homotopies.

Framing and coefficient
quantization change the equations themselves. On a stable rank-two
chart we construct regular classical and quantum Hamiltonian maps.
Their global extensions depend on which charges and which parameter
reductions are prescribed. Two obstruction theorems isolate those
dependencies. The remaining global comparison keeps its full scope.

\section{The derived commuting pair}

All polynomial algebras in this chapter are over $\mathbb C$.
Differentials raise cohomological degree by one.

Fix $N\geq1$,
put $G=GL_N$, and let $X,Y$ be matrices of independent commuting
coordinates. Matrix substitution in
\eqref{mht:eq:relation-resolution} gives
\begin{equation}\label{mht:eq:matrix-koszul}
 \mathcal K_N
 =\mathbb C[X_{ij},Y_{ij}]\otimes\Lambda(\Psi_{ij}),
 \qquad |\Psi_{ij}|=-1,\qquad
 QX=QY=0,\quad Q\Psi=[X,Y].
\end{equation}
The entries of $\Psi$ anticommute. The group acts by simultaneous
conjugation on all three matrices. Give $X,Y$ weight one and
$\Psi$ weight two, so that $Q$ preserves weight.

Let $I_N^{\mathrm{cl}}$ be the ideal generated by the entries of
$[X,Y]$. Then $H^0(\mathcal K_N)=\mathbb C[X,Y]/I_N^{\mathrm{cl}}$.
The full complex is needed even before taking gauge cochains:
\[
 Q\operatorname{Tr}\Psi=0,\qquad
 [\operatorname{Tr}\Psi]\ne0\quad\text{in }H^{-1}(\mathcal K_N^G).
\]
Indeed, setting $X=Y=0$ is a chain map to $\Lambda(\Psi_{ij})$
with zero differential, and the trace has nonzero image. More
precisely, splitting $\Psi$ into its traceless and scalar parts
factors $\mathcal K_N$ as the traceless constraint complex
tensor $\Lambda(\operatorname{Tr}\Psi)$.

Infinitesimal gauge transformations are imposed by
$C^\bullet(\mathfrak{gl}_N,\mathcal K_N)$, the
Chevalley--Eilenberg complex. Introduce its matrix $c$ of odd
generators of degree one. With graded matrix commutators the
total differential is
\begin{equation}\label{mht:eq:matrix-classical-brst}
 QX=[X,c],\quad QY=[Y,c],\quad Qc=-c^2,\quad
 Q\Psi=[X,Y]-[c,\Psi].
\end{equation}
These are the full gauge and relation terms. In each fixed
weight, the complex is bounded and its two differentials give
\[
 E_2^{a,b}
 =H^a\bigl(\mathfrak{gl}_N;H^b(\mathcal K_N)\bigr)
 \Longrightarrow
 H^{a+b}C^\bullet(\mathfrak{gl}_N,\mathcal K_N).
\]

Exactness of rational $G$-invariants identifies
$H^b(\mathcal K_N^G)$ with $H^b(\mathcal K_N)^G$.
It does not erase the positive Lie-cochain degrees.

The same information appears in the tangent complex at a
commuting pair:
\begin{equation}\label{mht:eq:matrix-tangent}
 \underbrace{\mathfrak{gl}_N}_{-1}
 \xrightarrow{\,d_0\,}
 \underbrace{\operatorname{Mat}_N\oplus\operatorname{Mat}_N}_{0}
 \xrightarrow{\,d_1\,}
 \underbrace{\mathfrak{gl}_N}_{1},
 \quad
 \begin{cases}
 d_0(a)=([a,X],[a,Y]),\\
 d_1(U,V)=[U,Y]+[X,V].
 \end{cases}
\end{equation}
Its composite is $[a,[X,Y]]=0$. On the middle term use
\[
 \omega((U,V),(U',V'))=\operatorname{Tr}(UV'-VU'),
\]
and pair the outer terms by $\operatorname{Tr}(ab)$,
symmetrically under reversal. This is graded skew because their
degrees are odd. Trace invariance gives
$\omega(d_0a,(U,V))=\operatorname{Tr}(a\,d_1(U,V))$.
The pairings are perfect on each term.

Both the stabilizer and relation directions participate in this
degree-zero symplectic pairing.

\section{Traces after the relation}

Let $F=\mathbb C\langle q,p\rangle$ and let
$F_{\mathrm{cyc}}=F/[F,F]$ denote its quotient by the linear
span of commutators. A free cyclic class need not remain nonzero
in the matrix relation complex. For example,
\begin{equation}\label{mht:eq:free-trace-boundary}
 Q\operatorname{Tr}(\Psi YX)
   =\operatorname{Tr}(X^2Y^2-XYXY).
\end{equation}

The classes $[q^2p^2]$ and $[qpqp]$ are distinct cyclic words.
Their trace difference takes value $-1$ at
$X=E_{12}$, $Y=E_{21}$, also after embedding these matrices in
rank four. Thus the right side is a nonzero free invariant
polynomial which becomes a boundary. Stable rank does not remove
this relation.

Write $A_+=(q,p)\mathbb C[q,p]$, graded by total degree.
Write $x_i,y_i$ for the coordinates of a diagonal matrix pair.
For a monomial $f=q^ap^b$ of positive degree, define
\[
 p_f^{(N)}=\sum_{i=1}^N x_i^a y_i^b.
\]
Let $\operatorname{Sym}(A_+)$ be the symmetric algebra on
linear symbols $\tau_f$, with $\tau_f$ of the same weight as $f$.
Its multiplication is the product of trace symbols. Thus
$\tau_{fg}$ and $\tau_f\tau_g$ are different elements.

Give $\operatorname{Sym}(A_+)$ the coproduct
$\Delta^\oplus(\tau_f)=\tau_f\otimes1+1\otimes\tau_f$.
An element $z$ is called \emph{block-sum primitive} if
$\Delta^\oplus z=z\otimes1+1\otimes z$.

\begin{proposition}\label{mht:prop:commuting-traces}
Restriction to diagonal pairs gives
\[
 H^0(\mathcal K_N^G)
 \cong\mathbb C[x_1,y_1,\ldots,x_N,y_N]^{S_N}.
\]
The homomorphism
\[
 \operatorname{Sym}(A_+)\longrightarrow H^0(\mathcal K_N^G),
 \qquad \tau_{q^ap^b}\longmapsto[\operatorname{Tr}X^aY^b],
\]
is surjective and is injective in every weight $d\leq N$.
Its degreewise stable value is $\operatorname{Sym}(A_+)$.
In this stable algebra the block-sum primitive subspace is $A_+$.
\end{proposition}
\begin{proof}
Rational reductive exactness gives
$H^0(\mathcal K_N^G)=(\mathbb C[X,Y]/I_N^{\mathrm{cl}})^G$.
The invariant part of the moment ideal equals that of its radical,
and restriction of the reduced commuting invariants to diagonal
pairs is an isomorphism. These are
\href{https://arxiv.org/abs/math/0409262}
{Gan--Ginzburg, Theorem 1.2.1 and Section 2.8}.
Only the invariant ideal is asserted radical here.

We prove the remaining statements directly in the diagonal ring.
Expand a product $p_{f_1}^{(N)}\cdots p_{f_r}^{(N)}$ by assigning
a particle index to each factor. The terms with distinct indices
give the corresponding invariant monomial sum. All other terms
use fewer indices and replace some $f_i$ by their products.
Induction on $r$, followed by division by the integer
multiplicities of repeated monomials, expresses every invariant
monomial sum as a polynomial in the $p_f^{(N)}$.

For injectivity in weight $d$, take a relation and choose its
largest number $r$ of factors $\tau_f$. Since every $f$ has
positive degree, $r\leq d\leq N$. Terms occupying exactly $r$
distinct labeled pairs can only come from products with $r$
factors. Their exponent multiset determines that product, with a
nonzero integer multiplicity. Its coefficient must therefore
vanish. Descending induction on $r$ proves injectivity.

Block sum sends each $\tau_f$ to
$\tau_f\otimes1+1\otimes\tau_f$.
In characteristic zero, the bidegree $(1,r-1)$ part of the
reduced coproduct on symmetric degree $r\geq2$ is the polynomial
differential. It is nonzero on every nonzero homogeneous
polynomial. Constants contribute no primitives. Hence the
primitive subspace is precisely the span of the $\tau_f$.
\end{proof}

The primitive coordinate is thus the vector space of nonconstant
commutative Hamiltonians, after imposing the matrix relation.
The map from $\operatorname{Sym}((F_{\mathrm{cyc}})_{>0})$
first factors through abelianization to
$\operatorname{Sym}(A_+)$, and then through the fixed-rank
power-sum map. The first kernel contains the cyclic difference
$[q^2p^2]-[qpqp]$ in \eqref{mht:eq:free-trace-boundary}.
Only the second kernel vanishes in the stated stable range.

At rank two the first relations occur in weight three. For
example, Cayley--Hamilton and its polarization give
\begin{align*}
 \operatorname{Tr}X^3
 &=\tfrac32\operatorname{Tr}X\,\operatorname{Tr}X^2
   -\tfrac12(\operatorname{Tr}X)^3,\\
 \operatorname{Tr}X^2Y
 &=\operatorname{Tr}X\,\operatorname{Tr}XY
   +\tfrac12\operatorname{Tr}Y\,\operatorname{Tr}X^2
   -\tfrac12(\operatorname{Tr}X)^2\operatorname{Tr}Y.
\end{align*}
Interchanging $X,Y$ gives the other two cubic identities.
Thus these four single traces vanish in the indecomposable
quotient of the positive-degree rank-two invariant algebra.
Their stable symbols are nonzero. Neither this quotient nor
ordinary invariant cohomology describes the entire complex
\eqref{mht:eq:matrix-classical-brst}.

\section{Central rank and the quadratic actions}

Give matrix phase space the canonical bracket
\[
 \{X_{ij},Y_{kl}\}=\delta_{il}\delta_{jk}.
\]
The full classical gauge algebra also pairs $c$ with $\Psi$,
symmetrically on these odd variables, as in
\eqref{mht:eq:matrix-classical-brst}. For a plane polynomial $f$
put $J_N(f)=\operatorname{Tr}(\sigma f)(X,Y)$, where $\sigma$
is the word average of Chapter~\ref{mht:ch:coherent-actions}.
These are closed Hamiltonian functions. They retain
\[
 J_N(1)=N,\qquad
 \{\operatorname{Tr}X,\operatorname{Tr}Y\}=N.
\]

Evaluation of all variables at zero is a cochain map to
$\mathbb C$, so $N$ is not exact. Consequently the first
components $q\mapsto\operatorname{Tr}X$ and
$p\mapsto\operatorname{Tr}Y$ cannot define a homotopy moment map
from $\mathbb C[q,p]/\mathbb C1$: their source bracket has
become zero. The central extension retaining $1$, or the pointed
source $(q,p)^2$, avoids this particular obstruction.

The scalar polynomials of degree at most two form a six-dimensional
Lie algebra. Their trace assignment is strict:
\[
 1\mapsto N,\quad q\mapsto\operatorname{Tr}X,\quad
 p\mapsto\operatorname{Tr}Y,\quad
 \frac{q^2}{2}\mapsto\frac{\operatorname{Tr}X^2}{2},\quad
 qp\mapsto\operatorname{Tr}XY,\quad
 \frac{p^2}{2}\mapsto\frac{\operatorname{Tr}Y^2}{2}.
\]
Differentiating these six functions verifies their brackets.
For example, the two pure quadratics bracket to
$\operatorname{Tr}XY$, and this function acts with weights
$-1,1$ on the two linear traces.

Allowing a matrix coefficient inside the trace gives a larger,
generally noninvariant space, and changes the closure question.
At $N=2$ let $\mathcal S$ be spanned by $1$ and
\[
 \operatorname{Tr}(AX),\quad\operatorname{Tr}(AY),\quad
 \operatorname{Tr}(AX^2),\quad\operatorname{Tr}(AY^2),\quad
 \tfrac12\operatorname{Tr}(A(XY+YX)),
 \qquad A\in\operatorname{Mat}_2.
\]

For $A=\operatorname{diag}(1,-1)$ one obtains
\[
 \left\{\tfrac12\operatorname{Tr}(AX^2),
             \tfrac12\operatorname{Tr}(AY^2)\right\}
       =X_{11}Y_{11}+X_{22}Y_{22}.
\]
This is not in $\mathcal S$, even modulo $I_2^{\mathrm{cl}}$.
Indeed, bidegree forces a putative expression to be
$\operatorname{Tr}(C(XY+YX))/2$. Setting $Y=I_2$ forces
$C=I_2$. Setting $X=Y=\left(\begin{smallmatrix}0&1\\1&0\end{smallmatrix}\right)$
then gives zero on the left and two on the right, at a commuting
pair in both cases.

\begin{proposition}\label{mht:prop:matrix-quadratic-closure}
The Lie algebra generated by $\mathcal S$ is the full space of
polynomials of degree at most two in the eight matrix entries.
It has dimension $36+8+1=45$, with constants as its center.
For a constant symplectic bracket, Weyl quantization preserves
the normalized Lie bracket on every such quadratic space.
\end{proposition}
\begin{proof}
Write $\sigma_0=I_2$ and set
\[
 \sigma_1=\begin{pmatrix}0&1\\1&0\end{pmatrix},\qquad
 \sigma_2=\begin{pmatrix}0&-i\\i&0\end{pmatrix},\qquad
 \sigma_3=\begin{pmatrix}1&0\\0&-1\end{pmatrix}.
\]
These satisfy $\operatorname{Tr}(\sigma_i\sigma_j)=2\delta_{ij}$.
Set $X=\sum r_i\sigma_i$ and $Y=\frac12\sum s_i\sigma_i$,
so $\{r_i,s_j\}=\delta_{ij}$ for $0\leq i,j\leq3$.
The Jordan operators $U\mapsto(AU+UA)/2$ for
$A=I_2,\sigma_i$ have matrices $I_4$ and
$S_i=E_{0i}+E_{i0}$.

For symmetric $T$ and arbitrary $C$ put
$U(T)=r^tTr/2$, $V(T)=s^tTs/2$, $W(C)=r^tCs$.
The initial space contains these forms for
$T,C\in\operatorname{span}(I_4,S_1,S_2,S_3)$. Their brackets are
\begin{align*}
 \{U(T),V(T')\}&=W(TT'),\\
 \{W(C),U(T)\}&=-U(CT+TC^t),\\
 \{W(C),V(T)\}&=V(C^tT+TC).
\end{align*}

For $i\ne j$, $S_iS_j+S_jS_i=E_{ij}+E_{ji}$, while
$S_i^2=E_{00}+E_{ii}$ and $\sum_iS_i^2-I_4=2E_{00}$.
The brackets therefore generate all symmetric $T$ in both
$U(T)$ and $V(T)$. Products of symmetric matrices span all
matrices, since $E_{ii}(E_{ij}+E_{ji})=E_{ij}$ for $i\ne j$.
They generate every $W(C)$.

The quadratic dimensions are $10+10+16$.
All eight linear coordinates were already present,
and their brackets supply the constant. Commuting with all
linear coordinates forces every derivative to vanish, so the
center consists of constants.

For a constant Poisson tensor $\Pi$ and central parameter $h$,
the polynomial Moyal product over $\mathbb C[h]$ is
$m\exp(h\Pi/2)$. Its commutator contains only odd powers of
$\Pi$. On quadratic inputs, the terms of order at least three
vanish, leaving $h\{F,G\}$. Weyl ordering identifies this
product with operator multiplication, and therefore preserves
the bracket divided by $h$.
\end{proof}

Projection to quadratics does not extend this result to nonlinear
Hamiltonian actions. For
$H=\operatorname{Tr}X^3/6$, $U=\operatorname{Tr}Y^2/2$ and
$V=\operatorname{Tr}X$, let $D$ be $\{H,-\}$ followed by that
projection. Then
\[
 DU=DV=0,\qquad D\{U,V\}=-\tfrac12\operatorname{Tr}X^2.
\]
The discarded cubic $\{H,U\}=\operatorname{Tr}(YX^2)/2$
is exactly the term needed for the derivation identity.
The quadratic space is a Lie subalgebra, not a quotient of the
full Hamiltonian action.

\section{Quantizing matrix entries}

Let $k_h=\mathbb C[h]$. The matrix Weyl algebra $W_{N,h}$ is
generated over $k_h$ by $X_{ij},Y_{ij}$ with
\begin{equation}\label{mht:eq:entrywise-weyl}
 [X_{ij},Y_{kl}]=h\delta_{il}\delta_{jk},\qquad
 [X_{ij},X_{kl}]=[Y_{ij},Y_{kl}]=0.
\end{equation}
The ordered monomials, with all $X$ before all $Y$, span the algebra.
Their differential-operator images
$X_{ij}=x_{ij}$, $Y_{ij}=-h\partial_{x_{ji}}$ are independent
after inverting $h$. Thus they form a basis.
The matrix Weyl algebra is $h$-torsion
free before any localization.

The normal-ordered moment matrix is
\begin{equation}\label{mht:eq:normal-moment}
 M_h=YX-XY+hNI_N,\qquad \operatorname{Tr}M_h=0.
\end{equation}
For $a\in\mathfrak g=\mathfrak{gl}_N$, put
$\ell_a=-\operatorname{Tr}(aM_h)$. In the differential
representation, $\ell_a=h\tau(a)$, where
$\tau(a)X=[X,a]$. Hence
$[\ell_a,\ell_b]=h\ell_{[a,b]}$.
Scalar conjugation has zero moment.

Define the left ideal and the ordinary invariant reduction by
\[
 I_{N,h}=\sum_{i,j}W_{N,h}(M_h)_{ij},\qquad
 \mathcal A_{N,h}=W_{N,h}^G/(I_{N,h}\cap W_{N,h}^G).
\]
The intersection is two-sided in the invariant algebra, because
invariants commute with all $\ell_a$. We do not impose the moment
as a two-sided ideal of the full Weyl algebra.

Let $D(\mathfrak g)$ denote polynomial differential operators
on $\mathfrak g$. Let $\mathfrak t\subset\mathfrak g$ be the
diagonal subspace.
The order filtration $F_mD(\mathfrak g)$ bounds the number of
derivatives and permits arbitrary polynomial coefficients.
The assignment above identifies $W_{N,h}$ with
$\bigoplus_{m\geq0}h^mF_mD(\mathfrak g)$.
Let $J_D=(D(\mathfrak g)\tau(\mathfrak g))^G$.

On the regular Cartan, ordinary radial restriction acts on
invariant functions. Put
$\delta=\prod_{i<j}(x_i-x_j)$ and
$\Phi(D)=\delta\operatorname{rad}(D)\delta^{-1}$.
The filtered Harish--Chandra theorem states that
\begin{equation}\label{mht:eq:filtered-hc}
 D(\mathfrak g)^G/J_D\xrightarrow{\ \Phi\ }D(\mathfrak t)^{S_N}
 \quad\text{is a strict filtered isomorphism},\qquad
 \operatorname{gr}J_D=(\sqrt{I_N^{\mathrm{cl}}})^G.
\end{equation}
Its associated graded map is diagonal restriction
(\href{https://arxiv.org/abs/math/0011114}
{Etingof--Ginzburg, Theorems 6.8--6.10}).
The invariant-ideal equality used in
Proposition~\ref{mht:prop:commuting-traces} removes the radical
in this formula.

An ideal is $h$-saturated if $hu$ belonging to it implies
that $u$ belongs to it.

\begin{theorem}\label{mht:thm:integral-matrix-reduction}
Conjugated radial restriction gives an isomorphism
\[
 \Phi_{N,h}:\mathcal A_{N,h}
 \xrightarrow{\ \sim\ }
 \bigl(k_h\langle x_i,\pi_i\rangle_{i=1}^N\bigr)^{S_N},
 \qquad [x_i,\pi_j]=h\delta_{ij},
\]
with all other pair commutators zero. The invariant moment ideal
is $h$-saturated, the quotient is free over $k_h$, and its special
fiber is $\mathbb C[x_i,\pi_i]^{S_N}$.
\end{theorem}
\begin{proof}
We first control coefficient orders in $J_D$. If
$D\in J_D\cap F_mD(\mathfrak g)$, its leading symbol belongs to
$(I_N^{\mathrm{cl}})^G$. Write it as
$\sum_a p_a\sigma_1(\tau(a))$, with each $p_a$ homogeneous of
momentum degree $m-1$. This is a finite expression.

Exactness
of rational invariants chooses an invariant coefficient vector
mapping to that symbol. Normal quantization is equivariant,
so it gives invariant $\sum_aP_a\tau(a)$ with the same symbol
and $\operatorname{ord}P_a\leq m-1$. Subtract and repeat at
lower order.

There is no order-zero remainder in the moment
ideal. We have proved a finite expression
$D=\sum_aP_a\tau(a)$ with those coefficient bounds.

Taking invariants in the Rees algebra is degreewise. Every
element of $I_{N,h}\cap W_{N,h}^G$ has its differential-operator
coefficient in $J_D$. Conversely, the expression just proved gives
\[
 h^mD=\sum_a(h^{m-1}P_a)(h\tau(a))\in I_{N,h}.
\]
Thus $I_{N,h}\cap W_{N,h}^G=\operatorname{Rees}_FJ_D$.
Strictness of \eqref{mht:eq:filtered-hc} identifies the quotient
with the Cartan invariant Rees algebra in the statement.

Cartan Weyl symmetrization is an $S_N$-equivariant $k_h$-linear
isomorphism from the ordinary polynomial module. A complex
basis of its invariant symbols gives a $k_h$-basis of the
quotient. Freeness implies $h$-saturation of the kernel:
if $hu$ maps to zero, then so does $u$.
Taking the intrinsic Rees special fiber proves the last assertion.
\end{proof}

In that special fiber, $\pi_i$ remains a momentum symbol.
Setting $h=0$ in the external representation
$\pi_i=-h\partial_{x_i}$ would instead kill it. Specializing
that representation does not compute the Rees special fiber.

\section{The radial trace and its Moyal bracket}

Use the same word average to define $J_N$ in $W_{N,h}$.
For $q^ap^b$, it averages traces of distinct words with $a$
factors $X$ and $b$ factors $Y$.
Thus $J_N(1)=N$.

Let $\operatorname{Op}_W$ denote the
analogous scalar Weyl ordering in one pair. On $k_h[q,p]$ put
\[
 f\star_hg=m\exp\!\left(\frac h2
 (\partial_q\otimes\partial_p-\partial_p\otimes\partial_q)\right)
 (f\otimes g),\qquad
 B_h(f,g)=\frac{f\star_hg-g\star_hf}{h}.
\]
The exponential terminates on polynomials, and the numerator
is divisible by $h$. This notation inverts no parameter.

\begin{theorem}\label{mht:thm:radial-trace}
For every $N$ and all polynomial $f,g$,
\[
 \Phi_{N,h}([J_N(f)])
       =\sum_i\operatorname{Op}_Wf(x_i,\pi_i),\qquad
 h^{-1}[J_N(f),J_N(g)]-J_N(B_h(f,g))
       \in I_{N,h}\cap W_{N,h}^G.
\]
The induced map
$k_h[q,p]/k_h1\longrightarrow\mathcal A_{N,h}/k_h1$
is injective as a map of Lie algebras with normalized brackets.
\end{theorem}
\begin{proof}
Pure coordinate traces restrict to $\sum_i x_i^d$.
The mixed second variation of an eigenvalue in entries
$(i,j),(j,i)$ is $\pm(x_i-x_j)^{-1}$, obtained by
differentiating $(z-x_i)(z-x_j)-uv=0$. Therefore
\[
 \operatorname{rad}(\operatorname{Tr}Y^2)
 =h^2\left(\sum_i\partial_i^2+
       2\sum_{i<j}\frac{\partial_i-\partial_j}{x_i-x_j}\right).
\]
The Vandermonde $\delta$ is harmonic: its Laplacian is an
alternating polynomial of smaller degree, hence is divisible
by $\delta$ and must be zero. Conjugating this radial operator
by $\delta$ gives $\Phi_{N,h}(\operatorname{Tr}Y^2)=\sum_i\pi_i^2$.

The derivation $D=h^{-1}[-\operatorname{Tr}Y^2/2,-]$
satisfies $DX=Y$, $DY=0$. Its exponential terminates on each
polynomial. Multiplicativity of $\Phi_{N,h}$ and torsion
freeness intertwine $D$ with $x_i\mapsto\pi_i$, $\pi_i\mapsto0$.
Apply this exponential to $\operatorname{Tr}X^{a+b}$ and
extract its coefficient of $t^b$. The same binomial factor
on both sides gives the claimed mixed trace formula.

Distinct particle pairs commute. In one pair, Weyl ordering
intertwines the Moyal product with operator multiplication.
The displayed defect thus has zero radial image, and
Theorem~\ref{mht:thm:integral-matrix-reduction} puts it in the
actual invariant moment ideal.

If the image of $J_N(f)$ is scalar, inverse Weyl symmetrization
gives $\sum_i f(x_i,\pi_i)=a\in k_h$. Set all pairs except
the first to zero. Then $f(q,p)+(N-1)f(0,0)=a$, so $f$
is constant. This proves injectivity of the scalar Lie quotient.
\end{proof}

This proves ideal membership in every bidegree. It does not
choose coefficients in a prescribed finite span of matrix
words, or compatible primitives in every BRST arity.

All these ordinary maps extend to $h$-adic completion, with a
polynomial coefficient at each order and no uniform degree bound.
To see the kernel assertion, put $J=I_{N,h}\cap W_{N,h}^G$.
Saturation gives $J\cap h^mW_{N,h}^G=h^mJ$ for every $m$.
The short exact sequence remains exact modulo $h^m$, and
$J/h^{m+1}J\to J/h^mJ$ is onto. Successively correcting lifts
in these kernels proves exactness of the inverse limit.
The completed kernel is therefore the closure of $J$.

\section{Rank splitting and the collective algebra}

The Cartan description separates particles. Set
$\mathcal A_{0,h}=k_h$ and $J_0(f)=0$.
Inclusion of $S_{N+M}$-invariants into
$S_N\times S_M$-invariants, transported through $\Phi$, gives
injective unital maps
\[
 \Delta^{\mathrm{red}}_{N,M}:\mathcal A_{N+M,h}
 \longrightarrow\mathcal A_{N,h}\otimes_{k_h}\mathcal A_{M,h}.
\]

Finite-group averaging identifies the target invariants with
this tensor product. Both ways of splitting three blocks are
the same subgroup inclusion, so the maps are coassociative
across ranks. The trace formula gives
\[
 \Delta^{\mathrm{red}}_{N,M}[J_{N+M}(f)]
     =[J_N(f)]\otimes1+1\otimes[J_M(f)].
\]
In particular the empty trace splits additively as $N+M$.
These maps do not define a coproduct on one fixed-rank algebra.

To retain products of one-particle sums, put $L=k_h[q,p]$ and
introduce a linear symbol $e(f)$ for each $f\in L$. Define
\begin{equation}\label{mht:eq:collective-moyal}
 \mathcal U_h=
 T_{k_h}(L)\big/
 \bigl(e(f)e(g)-e(g)e(f)-h\,e(B_h(f,g))\bigr).
\end{equation}
The generator $\nu=e(1)$ is central and is distinct from the
multiplicative unit.

\begin{theorem}\label{mht:thm:collective-moyal}
Ordered products of the monomial generators $e(q^ap^b)$ form
a $k_h$-basis of $\mathcal U_h$. It is a Hopf algebra with
\[
 \Delta e(f)=e(f)\otimes1+1\otimes e(f),\qquad
 \epsilon e(f)=0,\qquad S e(f)=-e(f).
\]
The maps $j_N:e(f)\mapsto[J_N(f)]$ are surjective algebra maps
to $\mathcal A_{N,h}$ and satisfy
\[
 \bigcap_{N\geq1}\ker j_N=0,\qquad
 (j_N\otimes j_M)\Delta
       =\Delta^{\mathrm{red}}_{N,M}j_{N+M}.
\]
These statements of joint injectivity and rank compatibility
also hold after $h$-adic completion, with completed tensor products.
\end{theorem}
\begin{proof}
The bracket $hB_h$ satisfies Jacobi because it is the Moyal
commutator. Order the monomial basis of $L$. Replacing an
inverted adjacent pair by the ordered pair and its bracket
decreases word length or, at fixed length, inversion number.
Disjoint replacements commute, and the difference for an
overlapping triple is Jacobi. These terminating reductions
give unique normal forms. Their leading coefficients are one,
so this proves the basis over $k_h$ without parameter inversion.

Primitive coproducts preserve the defining commutator relation:
the cross terms commute and the remaining terms are its copies
in the two tensor factors. Counit and antipode preserve it as
well. The Hopf identities hold on generators and, by induction,
on their products.

Theorem~\ref{mht:thm:radial-trace} gives the maps $j_N$.
For surjectivity, filter the Cartan Weyl algebra by total
polynomial degree, putting $h$ in the coefficient ring.
Its invariant associated graded is generated by polarized
power sums, by the distinct-index induction in
Proposition~\ref{mht:prop:commuting-traces}. Lift an invariant
leading symbol by a polynomial in the trace operators and
subtract it. Repeating at lower total degree proves generation.

Modulo $h$, the maps have source
$\operatorname{Sym}_{\mathbb C}(\mathbb C[q,p])$ and send
$e(f)$ to $p_f^{(N)}$, including $\nu\mapsto N$.
They are jointly injective. Indeed, group any proposed relation
by the largest number $r$ of nonconstant symbols, with
coefficients polynomial in $\nu$. For $N\geq r$, terms
occupying exactly $r$ distinct pairs identify each coefficient
polynomial at $\nu=N$, up to a nonzero integer. Vanishing for
all such $N$ makes every coefficient polynomial zero.
Descend in $r$.

For a nonzero element of $\mathcal U_h$, factor the common
lowest power of $h$ from its PBW coefficients. Torsion
freeness of every target would force its nonzero class modulo
$h$ to vanish in all classical targets, a contradiction.
A nonzero completed series also has a first nonzero coefficient,
which is a finite polynomial in PBW generators, so the same
argument applies. Finally the two rank maps in the statement
agree on $e(f)$ by splitting its one-particle sum. They agree
on the entire algebra.
\end{proof}

No finite collection of ranks is faithful: for a finite set
$T$ of ranks, the nonzero central polynomial
$\prod_{n\in T}(\nu-n)$ vanishes in all those representations.
Moreover $\nu$ is primitive, whereas $\Delta1=1\otimes1$.
The rank-$N$ quotient for $N>0$ is therefore not a Hopf quotient
with the same coproduct and counit.

Products cannot be removed by a linear projection which
preserves all quantum constants. Such a projection, if it killed
both $J_N(q)J_N(p)$ and $J_N(p)J_N(q)$, would send their
commutator $hN1$ to zero while fixing it. The collective
algebra retains these products instead.

\section{Quantum constraints and the full BRST algebra}

The ordinary quotient is the degree-zero part of a constraint
complex, but the complex has further cohomology. For
$\mathfrak g=\mathfrak{gl}_N$ define
$C^{-r}=W_{N,h}\otimes\Lambda^r\mathfrak g$, with
\begin{align*}
 d(w\otimes a_1\wedge\cdots\wedge a_r)
 ={}&\sum_i(-1)^{i+1}w\ell_{a_i}\otimes a_{\widehat i}\\
 &+h\sum_{i<j}(-1)^{i+j}
   w\otimes[a_i,a_j]\wedge a_{\widehat i\widehat j}.
\end{align*}

The omitted entries retain their order. This is the Chevalley
chain differential for the right action $w\cdot a=w\ell_a$
and bracket $h[-,-]$. In degree two its square is
$w[\ell_a,\ell_b]-hw\ell_{[a,b]}=0$. In higher degrees the
remaining overlaps are the same identity and Jacobi. Hence
$d^2=0$, and reductive exactness gives
$H^0(C^G)=\mathcal A_{N,h}$.

Splitting $\mathfrak g=\mathfrak{sl}_N\oplus\mathbb CI_N$
factors this complex as
\[
 C(\mathfrak g)=C(\mathfrak{sl}_N)\otimes\Lambda(\theta_0),
 \qquad |\theta_0|=-1,\quad d\theta_0=0.
\]
The scalar comoment is zero. Thus $H^{-1}(C^G)$ has the
direct summand $\mathcal A_{N,h}\theta_0$.
Exterior multiplication with multiplication of coefficients
does not make $C$ a differential algebra: its differential
contains right multiplication by the constraints.

To obtain a differential algebra, choose a basis $a_i$ with
$[a_i,a_j]=f_{ij}^{\ k}a_k$ and write $\ell_i=\ell_{a_i}$.
Adjoin odd ghosts $c^i,b_i$ of degrees $1,-1$.
Their $k_h$-algebra, denoted $\operatorname{Cl}_h(c,b)$, has relations
\[
 c^ic^j=-c^jc^i,\qquad b_ib_j=-b_jb_i,\qquad
 b_ic^j+c^jb_i=h\delta_i^j.
\]

Set $\mathcal B_{N,h}=W_{N,h}\otimes_{k_h}\operatorname{Cl}_h(c,b)$,
with the graded tensor product. Repeated indices are summed in
\begin{equation}\label{mht:eq:quantum-brst-charge}
 \Theta=\sum_i c^i\ell_i-\tfrac12f_{ij}^{\ k}c^ic^jb_k,
 \qquad Q=h^{-1}[\Theta,-]_{\mathrm{gr}}.
\end{equation}

Normal exterior words in $c,b$ form a free $k_h$-basis:
represent $c^i$ by exterior multiplication and $b_i$ by
$h\partial_{c^i}$. For one pair the four normal words give
independent two-by-two matrices after inverting $h$.
Tensoring proves independence for all pairs.
The relations give spanning before inversion. Consequently
$\mathcal B_{N,h}$ is $h$-torsion free. Modulo $h$ it is
supercommutative, so the divided commutator defining $Q$
is integral.

\begin{proposition}\label{mht:prop:quantum-brst}
The charge satisfies $\Theta^2=0$, and $Q$ is a
degree-one derivation with $Q^2=0$. At rank one the full
$\mathfrak{gl}_1$ model has $Q=0$ and its degree-zero ghost
factor is free on $1,u=cb$ with $u^2=hu$.
\end{proposition}
\begin{proof}
Use the faithful ghost representation above and the left regular
matter representation. Then
$\Theta=\sum c^i\ell_i+h\,d_{\mathrm{CE}}$.
Jacobi gives $d_{\mathrm{CE}}^2=0$. The square of the first
term is $\frac h2f_{ij}^{\ k}c^ic^j\ell_k$, and the mixed
term is its negative. Thus $\Theta^2=0$ in a faithful
representation and hence in the algebra.

The inner graded commutator is a derivation, and its square is
$[\Theta^2,-]$. Torsion freeness gives $Q^2=0$.

At rank one both the moment and the structure constants vanish.
The ghost relation gives
$cbcb=c(h-cb)b=hcb$, with $1,cb$ independent in the normal basis.
\end{proof}

Already at rank one, $H^0(\mathcal B_{N,h},Q)$ therefore
differs from ordinary reduction. Choosing the effective
traceless gauge algebra removes this particular scalar ghost
factor by changing the gauge input.

The cubic trace defect has a primitive in the full model.
Use a trace-dual pair of bases $a_i,a^i$, raise moment indices
with that pairing, and put
\[
 T_N=\sum_i b_i\ell^i,\qquad C_N=\sum_i\ell_i\ell^i.
\]
From \eqref{mht:eq:quantum-brst-charge},
$Qb_i=\ell_i-f_{ij}^{\ k}c^jb_k$ and
$Q\ell^i=c^j[\ell_j,\ell^i]/h$. The graded Leibniz rule,
invariance of the pairing, and
$\sum_i[\ell_i,\ell^i]=0$ give
\begin{equation}\label{mht:eq:casimir-primitive}
 QT_N=C_N.
\end{equation}
Indeed the two terms containing $c^jb_k$ cancel after moving
$b_i$ past $c^j$. The resulting contraction is the vanishing
sum of moment commutators.

\begin{proposition}\label{mht:prop:cubic-brst}
For every rank, with the Moyal source bracket,
\[
 h^{-1}[J_N(q^3),J_N(p^3)]-J_N(B_h(q^3,p^3))
       =Q\bigl(-\tfrac32T_N\bigr).
\]
\end{proposition}
\begin{proof}
Let $a=\operatorname{Tr}(X^2Y^2)$ and
$b=\operatorname{Tr}(XYXY)$ denote commuting entry symbols.
Write $\operatorname{Op}^{\mathrm{ent}}_W$ for entrywise
Weyl quantization. The entries of $X+sY$ commute with one
another, so polarization proves that the matrix-word average
equals the entrywise quantization of its commutative symbol.
In particular $J_N(q^2p^2)=
\operatorname{Op}^{\mathrm{ent}}_W(2a/3+b/3)$.

The first entrywise contraction of the two cubic traces is
$9a$. Contracting their third derivatives gives
$18N^3+18N$. Each third derivative has two cyclic index
identifications, with multiplicity three.
The normalized third-order coefficient is $h^2/24$.
Thus
\[
 h^{-1}[\operatorname{Tr}X^3,\operatorname{Tr}Y^3]
 =9\operatorname{Op}^{\mathrm{ent}}_W(a)
      +\tfrac34h^2(N^3+N).
\]

For the paired quadratic moment symbols the first contraction
cancels. The second is
\[
 -\frac{h^2}{4}\sum_i\operatorname{Tr}
              (\operatorname{ad}a_i\operatorname{ad}a^i).
\]
This trace is on $\mathfrak g$. Writing each adjoint as left
minus right multiplication gives the sum $2N(N^2-1)$, and hence
\[
 C_N=\operatorname{Op}^{\mathrm{ent}}_W
       \operatorname{Tr}([X,Y]_{\mathrm{cl}}^2)
             -\tfrac12h^2(N^3-N).
\]

Finally,
\[
 3(a-b)=-\tfrac32\operatorname{Tr}([X,Y]_{\mathrm{cl}}^2).
\]
Subtracting $J_N(9q^2p^2)$ therefore leaves
$-\frac32 C_N+\frac32h^2N$.
But $B_h(q^3,p^3)=9q^2p^2+\frac32h^2$, and $J_N(1)=N$.
The scalar terms cancel. Equation~\eqref{mht:eq:casimir-primitive}
now gives the primitive.
\end{proof}

The choice of source bracket is essential. At rank one,
\[
 B_h(q^3,p^4)-\{q^3,p^4\}=6h^2p,
\]
and $Q=0$ in the full BRST model. The unchanged Weyl-symbol
map from the undeformed Poisson source therefore has a nonzero
binary defect which no homotopy can remove. This does not
exclude a changed first map or a different recipient.

\section{A formal slice with stabilizer ghosts}

Fix a labeled diagonal matrix with distinct entries $x_i^0$
and complete the matrix base at that point. Formal eigenvectors
with fixed normalizations give coordinates
\[
 X=g(s)D(x)g(s)^{-1},\qquad D(x)=\operatorname{diag}(x_i).
\]
There are $N^2-N$ orbit coordinates $s_j$. To construct them,
solve the eigenvector equations order by order.
The off-diagonal linear terms divide by the nonzero constants
$x_i^0-x_k^0$. This recursive solution also proves
invertibility of the formal coordinate change.

On the super base $(X,c)$, the differential is
$QX=[X,c]$, $Qc=-c^2$. Put
\[
 \gamma_j=Qs_j,\qquad c'=g^{-1}cg+g^{-1}Qg.
\]
Differentiating the diagonalization gives $QD=[D,c']$.
Its diagonal part is $Qx=0$ and its off-diagonal part forces
$c'$ to be diagonal.

The ghost change from $c$ to
$(\gamma_j,c'_i)$ is invertible: its linearization separates
the orbit directions from the diagonal stabilizer. The transformed
Maurer--Cartan equation gives $Qc'=-(c')^2=0$.
Also $Q\gamma=0$. Thus the vector field on this base is
$\sum_j\gamma_j\partial_{s_j}$.

Take the order Rees algebra of differential operators with
formal base coefficients, and complete it in $h$. Its symbols
have polynomial momentum degree at each parameter order.
The chain rule induces an exact algebra isomorphism, with
$Y=-h\partial_X$ and
$b=h\partial_c$. Matter momenta include the ghost terms
required by the change of $c$. Changing $X$ alone would not
produce this isomorphism.

With
$\pi_j=-h\partial_{s_j}$, $\beta_j=h\partial_{\gamma_j}$,
the charge becomes $\Theta=-\sum_j\gamma_j\pi_j$.
The orbit factors have
\[
 Qs_j=\gamma_j,\quad Q\beta_j=-\pi_j,\quad
 Q\gamma_j=Q\pi_j=0,\qquad
 [s_j,\pi_k]=h\delta_{jk},\quad
 \beta_j\gamma_k+\gamma_k\beta_j=h\delta_{jk}.
\]

The residual factor is
\[
 \mathcal R_{N,h}
 =W_h(T^*\mathfrak t)^\wedge_{x^0}
   \widehat\otimes
   \operatorname{Cl}_h(c'_1,\ldots,c'_N,b'_1,\ldots,b'_N),
 \qquad Q|_{\mathcal R_{N,h}}=0.
\]
Here the first factor has formal eigenvalue coefficients and
polynomial momenta at each parameter order. The second factor
retains the diagonal stabilizer ghosts and their conjugates.
The residual momenta are $p_i=-h\partial_{x_i}$ and
$b'_i=h\partial_{c'_i}$.

\begin{proposition}\label{mht:prop:formal-brst-slice}
The residual inclusion is a strict differential-algebra
quasi-isomorphism on this labeled formal chart. On it,
\[
 f\longmapsto\sum_i\operatorname{Op}_Wf(x_i,p_i)
\]
is a strict Lie map from the plane Moyal algebra, with the
normalized operator bracket.
\end{proposition}
\begin{proof}
Pass to Weyl symbols in the orbit variables. Let $\mathsf N$
count all four kinds of orbit generators and define on positive
count
\[
 H=\mathsf N^{-1}\sum_j
       (s_j\partial_{\gamma_j}-\beta_j\partial_{\pi_j}).
\]
The anticommutator of $Q$ with the numerator is $\mathsf N$.
Hence $QH+HQ=1-ip$, where $p$ extracts orbit count zero
and $i$ includes the residual factor. This calculation is
coefficientwise in formal $s$ and $h$, keeps finite momentum
degree at each order, and inverts only positive integers.
The separate residual factor makes $i$ a strict algebra map.

Finally, distinct particle pairs commute and each individual
pair realizes the plane Moyal product. Their sum gives the
stated Lie identity.
\end{proof}

The projection $p$ in this proof is linear, not multiplicative:
annihilating both $s_j$ and $\pi_j$ cannot preserve
$[s_j,\pi_j]=h$. The action is defined on a labeled formal
chart. Extending it across the discriminant, with comparisons
between charts and with prescribed global matrix charges,
requires compatible BRST homotopies in every arity. The
global ordinary reduction of
Theorem~\ref{mht:thm:integral-matrix-reduction} supplies a
different, already defined recipient.

\section{Representations of the coefficient Weyl algebra}

Let $\varepsilon$ denote the parameter in the coefficient
relation $qp-pq=\varepsilon$. It is independent of the matrix
quantization parameter $h$. Its derived $N$-dimensional
representation algebra has commuting coordinate entries and
differential
\[
 \mathcal K_{N,\varepsilon}
 =\mathbb C[\varepsilon,X,Y]\otimes\Lambda(\Psi),\qquad
 Q\Psi=[X,Y]-\varepsilon I_N.
\]

The element $v=-\operatorname{Tr}\Psi/N$ satisfies
$Qv=\varepsilon$. Consequently
$Q(va)+vQa=\varepsilon a$ for every homogeneous $a$.
All its cohomology is killed by $\varepsilon$, and after
inverting $\varepsilon$, multiplication by $v/\varepsilon$
contracts the complex. This is the homological form of
the finite-dimensional trace obstruction to the Weyl relation.
It does not apply to the entrywise relations
\eqref{mht:eq:entrywise-weyl}.

An infinite-dimensional recipient survives integrally.
Complete the coefficient Weyl algebra in $\varepsilon$ and
write it as $\widehat W_\varepsilon$. On
$\mathcal P=\mathbb C[p][[\varepsilon]]$ let
$p$ act by multiplication and $q$ by
$\varepsilon\partial_p$. Use complete modules, continuous maps,
strict exact sequences and resolutions by topologically free
modules. The PBW basis gives the strict resolution
\[
 0\longrightarrow\widehat W_\varepsilon
 \xrightarrow{\,R_q\,}\widehat W_\varepsilon
 \longrightarrow\mathcal P\longrightarrow0.
\]

Right multiplication by $q$ is injective with closed image:
in the normal basis $p^aq^b$ it increases $b$ by one.
The quotient consists of the $b=0$ terms. Continuous Hom into
$\mathcal P$ gives
$\mathcal P\xrightarrow{\varepsilon\partial_p}\mathcal P$.

Polynomial integration coefficient by coefficient is a
continuous right inverse to $\partial_p$. Therefore
\begin{equation}\label{mht:eq:integral-fock-ext}
 \operatorname{Ext}^0(\mathcal P,\mathcal P)=\mathbb C[[\varepsilon]],
 \qquad
 \operatorname{Ext}^1(\mathcal P,\mathcal P)=\mathbb C[p],
 \qquad \operatorname{Ext}^{i}=0\quad(i\ne0,1).
\end{equation}
The degree-one group is annihilated by $\varepsilon$ and
disappears after its inversion. It is part of the integral
module, not an additional generic-parameter state.

\section{A framed moment equation}

A single framing vector changes the finite-dimensional equation.
Fix $\zeta\in\mathbb C$, add a column $I$ and a row $J$, and impose
\begin{equation}\label{mht:eq:framed-moment}
 [X,Y]+IJ=\zeta I_N.
\end{equation}

Here $X,Y,I,J$ have commuting coordinate entries. Use
$\{X_{ij},Y_{kl}\}=\delta_{il}\delta_{jk}$,
$\{I_i,J_j\}=\delta_{ij}$, and zero for the other generator
brackets. A point is stable if the vectors $W(X,Y)I$, for
words $W$, span $\mathbb C^N$. The action is
\[
 g\cdot(X,Y,I,J)=(gXg^{-1},gYg^{-1},gI,Jg^{-1}).
\]
Let $\mathcal X_{N,\zeta}$ be the stable quotient.

This quotient is smooth and symplectic. Indeed, an element
fixing $X,Y,I$ fixes every spanning word vector, and is the
identity. The same argument gives injectivity of the
infinitesimal action. Its symplectic transpose is the moment
differential, which is therefore surjective.

The quotient has
dimension $(2N^2+2N)-2N^2=2N$. Choosing $N$ word vectors
which form a basis gives a gauge chart. Changes between these
bases give its regular transition maps.

On the ordered chart where $X$ has distinct eigenvalues,
diagonalize it. Stability forces every $I_i$ to be nonzero.
For if $I_i=0$, the off-diagonal moment equations give
$Y_{ij}=0$ for all $j\ne i$.
Every word applied to $I$ would then have zero $i$th component.

Diagonal gauge sets $I_i=1$,
and the equation becomes
\begin{equation}\label{mht:eq:framed-chart}
 X=\operatorname{diag}(x_i),\quad I_i=1,\quad J_i=\zeta,
 \qquad Y_{ii}=y_i,\quad
 Y_{ij}=-\frac{\zeta}{x_i-x_j}\quad(i\ne j).
\end{equation}
The ambient form $\operatorname{Tr}(dY\wedge dX)+dJ\wedge dI$
restricts to $\sum_i dy_i\wedge dx_i$, with
$\{x_i,y_j\}=\delta_{ij}$. Permutations give the unlabeled
spectral chart. Setting $\zeta=\varepsilon$ also gives a
framed recipient over the coefficient-parameter ring, on this
same discriminant-localized chart.

The necklace algebra, with its constant class retained, has a
global strict Hamiltonian map on $\mathcal X_{N,\zeta}$:
\begin{equation}\label{mht:eq:framed-necklace}
 [W]\longmapsto H_W=\operatorname{Tr}W(X,Y),\qquad H_1=N.
\end{equation}
Its necklace bracket is
$[\partial_q^{\mathrm{cyc}}W\,\partial_p^{\mathrm{cyc}}V
 -\partial_p^{\mathrm{cyc}}W\,\partial_q^{\mathrm{cyc}}V]$.
Cutting the two cyclic words at their differentiated factors
proves $\{H_W,H_V\}=H_{\{W,V\}_{\mathrm{neck}}}$.
The associated field has
$\delta X=-\partial_Y^{\mathrm{cyc}}W$,
$\delta Y=\partial_X^{\mathrm{cyc}}W$ and $\delta I=\delta J=0$.
The cyclic derivative identity makes $\delta[X,Y]=0$.
Thus it preserves the constraint and commutes with conjugation.

The full framed BFV algebra adds odd matrices $c,\Psi$ of
degrees $1,-1$, with
$\{c_{ij},\Psi_{kl}\}=\delta_{il}\delta_{jk}$. Its differential
is the bracket with
\[
 \Theta_{\mathrm{fr}}
 =\operatorname{Tr}\bigl(c([X,Y]+IJ-\zeta I_N)-\Psi c^2\bigr).
\]
Equivariance of the moment gives $Q^2=0$ as in the unframed
model. The Hamiltonians $H_W$ are closed: they are gauge
invariant, commute with the moment functions, and contain no
ghosts. No completeness of their polynomial flows is asserted.

\section{The all-rank defect of symmetrized traces}

For commuting entry symbols, average the traces of words
containing $a$ copies of $X$ and $b$ copies of $Y$.
Denote the result by $J_{ab}$, with $J_{00}=N$. The frame in
\eqref{mht:eq:framed-moment} changes which trace defects vanish.

\begin{proposition}\label{mht:prop:framed-residual}
On $\mathcal X_{N,\zeta}$,
\begin{align}
 \{J_{30},J_{03}\}-9J_{22}
       &=-\tfrac32N(N-1)\zeta^2,\label{mht:eq:framed-cubic}\\
 \{J_{30},J_{04}\}-12J_{23}
       &=-6(N-1)\zeta^2J_{01}.\label{mht:eq:framed-cubic-quartic}
\end{align}
For $N\geq2$, a plane Moyal bracket of parameter $\theta$
compatible with these unchanged trace images would have to
satisfy $\theta^2=-(N-1)\zeta^2$. With that value,
\begin{equation}\label{mht:eq:framed-quartic-residual}
 \{J_{40},J_{04}\}-16J_{33}-24\theta^2J_{11}
       =\frac{56}{5}\zeta^2(NJ_{11}-J_{10}J_{01}).
\end{equation}
For $\zeta\ne0$ this is nonzero at every rank $N\geq2$.
Hence no single Moyal parameter makes the unchanged map from
commutative polynomials to symmetrized traces a Lie map.
\end{proposition}
\begin{proof}
Put $C=[X,Y]$ and $U=IJ=\zeta I_N-C$.
Tracing the constraint gives $JI=N\zeta$, so $U^2=N\zeta U$.
Expansion of the six words of bidegree $(2,2)$ and
the ten of bidegree $(2,3)$ gives
\[
 \{J_{30},J_{03}\}-9J_{22}=-\tfrac32\operatorname{Tr}C^2,
 \qquad
 \{J_{30},J_{04}\}-12J_{23}=-6\operatorname{Tr}(YC^2).
\]
The rank-one expression for $U$ implies
\[
 C^2=(N-1)\zeta^2I_N-(N-2)\zeta C.
\]
Since $\operatorname{Tr}C=\operatorname{Tr}(YC)=0$, these
identities prove the two cubic formulas.

For the quartic pair use
$\operatorname{Tr}(UXUY)=\operatorname{Tr}(UX)\operatorname{Tr}(UY)$.
Together with
$\operatorname{Tr}(CX)=\operatorname{Tr}(CY)=
 \operatorname{Tr}(C(XY+YX))=0$, this gives
\begin{align*}
 \operatorname{Tr}(C^2(XY+YX))&=2(N-1)\zeta^2J_{11},\\
 \operatorname{Tr}(CXCY)&=\zeta^2(J_{10}J_{01}-J_{11}).
\end{align*}
The twenty words in $J_{33}$ divide into four cyclic classes
of sizes $6,6,6,2$, represented by
$X^3Y^3$, $X^2YXY^2$, $X^2Y^2XY$ and $(XY)^3$.
Expansion in these four classes yields
\[
 \operatorname{Tr}(X^3Y^3)-J_{33}
 =-\tfrac25\operatorname{Tr}(C^2(XY+YX))
  -\tfrac7{10}\operatorname{Tr}(CXCY).
\]
Since $\{J_{40},J_{04}\}=16\operatorname{Tr}(X^3Y^3)$,
substitution gives
\[
 \{J_{40},J_{04}\}-16J_{33}
 =\tfrac85\zeta^2\bigl((15-8N)J_{11}-7J_{10}J_{01}\bigr).
\]

The scalar Moyal formulas are
$B_\theta(q^3,p^3)=9q^2p^2+\frac32\theta^2$ and
$B_\theta(q^4,p^4)=16q^3p^3+24\theta^2qp$.
Because $J_{00}=N$, the cubic equation forces the stated
$\theta^2$. Subtracting the quartic Moyal term gives
\eqref{mht:eq:framed-quartic-residual}.
On \eqref{mht:eq:framed-chart}, choose $x_i=i$,
$y_1=1$ and $y_i=0$ for $i>1$. Then
$NJ_{11}-J_{10}J_{01}=-N(N-1)/2$, proving nonvanishing.
\end{proof}

The next unequal-degree pair records the same loss of a
product of traces. With the forced value of $\theta$,
\begin{equation}\label{mht:eq:framed-next-residual}
 \{J_{30},J_{05}\}-15J_{24}-15\theta^2J_{02}
          =6\zeta^2(NJ_{02}-J_{01}^2).
\end{equation}
Indeed, the fifteen words in $J_{24}$ have cyclic class sizes
$6,6,3$, and give
\[
 \operatorname{Tr}(X^2Y^4)-J_{24}
 =-\tfrac35\operatorname{Tr}(C^2Y^2)
  -\tfrac25\operatorname{Tr}(CYCY).
\]
The same rank-one identities evaluate these two traces as
$(N-1)\zeta^2J_{02}$ and $\zeta^2(J_{01}^2-J_{02})$.
Together with $B_\theta(q^3,p^5)=15q^2p^4+15\theta^2p^2$,
they prove \eqref{mht:eq:framed-next-residual}.

The quartic residual also excludes a homotopy moment map into
the framed BFV complex with these fixed first components.
Its binary identity would make the residual a boundary.
Evaluation at a constraint point, with ghosts and antighosts
zero, is a chain map to $\mathbb C$. At the point used in the
proof the residual has nonzero value. The obstruction concerns
these fixed first components.

\section{A particle map and a stable collision}

The distinct-eigenvalue chart has a strict map with a different
first component. Define spectral projectors and measurements by
\[
 P_i(X)=\prod_{j\ne i}\frac{X-x_jI_N}{x_i-x_j},
 \qquad y_i=\operatorname{Tr}(P_iY),\qquad
 \widehat J(f)=\sum_i f(x_i,y_i).
\]
The sum is conjugation invariant and invariant under permutations,
so it is a regular function on the spectral chart.

\begin{proposition}\label{mht:prop:framed-particle}
The map $\widehat J$ preserves the plane Poisson bracket and
sends $1$ to $N$. For the product Darboux Moyal structure of
parameter $h$ on that chart it also satisfies
\[
 h^{-1}[\widehat J(f),\widehat J(g)]_{\star_{\mathrm{chart}}}
       =\widehat J(B_h(f,g)).
\]
At $N=2$ and $\zeta\ne0$, its image of $p^2$ has no regular
extension to the entire stable quotient.
\end{proposition}
\begin{proof}
The Darboux coordinates of \eqref{mht:eq:framed-chart}
give $\{x_i,y_j\}=\delta_{ij}$. Different pairs Poisson-commute,
so the bracket of the two sums is
$\sum_i\{f,g\}(x_i,y_i)$. For the chart Moyal product use
the same exponential tensorwise on the discriminant-localized
coordinate ring, completed in $h$. Each coefficient is a
finite-order differential operator and respects localization
and permutations. Different pairs star-commute, proving the
quantum formula.

The change from matrix traces is already visible in degree two:
\[
 \widehat J(p^2)=\operatorname{Tr}Y^2
                +2\zeta^2\sum_{i<j}(x_i-x_j)^{-2}.
\]
To test regular extension, consider the family
\[
 X(a)=\begin{pmatrix}0&1\\a&0\end{pmatrix},\quad
 Y(a)=\begin{pmatrix}0&0\\\zeta&0\end{pmatrix},\quad
 I=\begin{pmatrix}0\\1\end{pmatrix},\quad J=(0,2\zeta).
\]

It satisfies \eqref{mht:eq:framed-moment} for every $a$.
The columns $I,XI$ have determinant $-1$, so the family remains
in one stable gauge chart even at $a=0$.

For $a=b^2\ne0$, the eigenvalues are $\pm b$, with projectors
$(I_N\pm X/b)/2$. Since $\operatorname{Tr}Y=0$ and
$\operatorname{Tr}(XY)=\zeta$, the two measurements are
$\pm\zeta/(2b)$. Therefore
$\widehat J(p^2)=\zeta^2/(2a)$, whereas $\operatorname{Tr}Y^2=0$.
A regular function on the full stable quotient would pull
back regularly at $a=0$. This pole excludes that extension.
\end{proof}

The particle sum is a Lie map. Products have terms on distinct
particles, so it is not an associative map from the one-particle
algebra. Its chart quantization has not been
identified with a global quantization of
$\mathcal X_{N,\zeta}$.

\section{Regular maps on the cyclic collision chart}
\label{mht:frgl:sec:local-chart}

The local classification works over any field $K$ of characteristic zero.
Fix $\zeta\in K$. Functions have cohomological degree zero. The bracket
on matrix entries is $\{X_{ij},Y_{kl}\}=\delta_{il}\delta_{jk}$ and the
framing bracket is $\{I_i,J_j\}=\delta_{ij}$. The quotient is the stable
rank-two quotient of $[X,Y]+IJ=\zeta\mathbf1$. Stability means that words
in $X,Y$ applied to $I$ span $K^2$.

Let $U_+$ be the open subset on which $I,XI$ are independent. Put
$s=\operatorname{Tr}X/2$. The basis $((X-s)I,I)$ gives the unique representative
\begin{equation}\label{mht:frgl:eq:local-gauge}
 X=s\mathbf1+\begin{pmatrix}0&1\\a&0\end{pmatrix},\qquad
 Y=t\mathbf1+\begin{pmatrix}0&k\\\zeta+ak&0\end{pmatrix},\qquad
 I=\binom01,\quad J=(0,2\zeta).
\end{equation}
Consequently $U_+=\operatorname{Spec}K[s,t,a,k]$.

To check this assertion,
Cayley--Hamilton first gives the displayed form of $X$. Write the traceless
part of $Y$ as $\left(\begin{smallmatrix}r&k\\\ell&-r\end{smallmatrix}\right)$.
The top-right moment equation gives $r=0$, the top-left gives
$\ell=\zeta+ak$, and the bottom row gives the displayed $J$. Conversely
every point in \eqref{mht:frgl:eq:local-gauge} satisfies the moment equation and
$\det(I,XI)=-1$.

The symplectic form pulls back to
\[
 \operatorname{Tr}(dY\wedge dX)+dJ\wedge dI
   =2\,dt\wedge ds+dk\wedge da.
\]
Thus $\{s,t\}=1/2$, $\{a,k\}=1$, and all other coordinate brackets
vanish.

Define centered invariants
\[
 b=\tfrac12\operatorname{Tr}((X-s)(Y-t))=\zeta/2+ak,
 \qquad c=\tfrac12\operatorname{Tr}(Y-t)^2=\zeta k+ak^2.
\]
Then
\begin{equation}\label{mht:frgl:eq:orbit}
 ac-b^2=-\zeta^2/4,\qquad
 \{a,b\}=a,\quad\{b,c\}=c,\quad\{a,c\}=2b.
\end{equation}

The collision of Proposition~\ref{mht:prop:framed-particle} is the line $s=t=k=0$. To recover its
measured entries explicitly, adjoin $r$ with $r^2=a$ and restrict to
$r\ne0$. Then
\[
 X=\begin{pmatrix}0&1\\a&0\end{pmatrix},\qquad
 Y=\begin{pmatrix}0&0\\\zeta&0\end{pmatrix},\qquad
 P_\pm=\tfrac12(\mathbf1\pm X/r).
\]
The spectral projectors obey
\[
 p_\pm=\operatorname{Tr}(P_\pm Y)=\pm\frac{\zeta}{2r},
 \qquad p_+^2+p_-^2=\frac{\zeta^2}{2a},
 \qquad \operatorname{Tr}Y^2=0.
\]
Thus the required measurements have their pole at a stable point:
$I,XI$ remain independent when $a=0$. The spectral square is a
different charge from the regular matrix square.

The symmetrized traces already defined above have $J_{00}=2$
in rank two. On the cyclic chart,
\[
 J_{10}=2s,\quad J_{01}=2t,\quad J_{20}=2s^2+2a,
 \quad J_{11}=2st+\zeta+2ak,\quad
 J_{02}=2t^2+2\zeta k+2ak^2.
\]

For a parameter $\tau\in K$, let $\mathfrak g_\tau$ be $K[q,p]$ with
the bracket
\begin{equation}\label{mht:frgl:eq:moyal-tau}
 \mathcal B_\tau(f,g)=\sum_{r\geq0}
 \frac{\tau^r}{2^{2r}(2r+1)!}\,
 m\left(\partial_q\otimes\partial_p-
       \partial_p\otimes\partial_q\right)^{2r+1}(f\otimes g).
\end{equation}

Here $m$ multiplies the two tensor factors. The sum is finite. Jacobi follows by adjoining $\theta$ with $\tau=\theta^2$ and taking the divided commutator of the associative Moyal product. The polynomial identity descends over $\mathbb Q[\tau]$. Thus $\mathcal B_\tau=B_\theta$ in the earlier notation when
$\tau=\theta^2$, and $\mathcal B_0$ is the Poisson bracket.

The target bracket in this
section is always the ordinary Poisson bracket, including when
$\tau\ne0$. No target star product is being identified with it.

Consider the five-dimensional Lie subalgebra
\[
 \mathfrak b=\langle 1,q,p,q^2,qp\rangle\subset\mathfrak g_\tau.
\]
Its bracket is independent of $\tau$. Suppose that a $K$-linear Lie map
$\Phi:\mathfrak g_\tau\to K[[s,t,a,k]]$ agrees with $J$ on
$\mathfrak b$. A germ of a regular algebraic map is a special case.

Set $F=\Phi(p^2)$. The source brackets with $q,p,q^2,qp$ give
\begin{align*}
 F_t&=4t,&F_s&=0,&F_k&=2\zeta+4ak,\\
 \{J_{11},F\}&=2F.
\end{align*}

The first three equations imply
$F=J_{02}+H(a)$. The last is $aH'(a)+H(a)=0$. In formal power series
the coefficient of $a^n$ is multiplied by $n+1$, so $H=0$.
On the punctured chart its only Laurent solution is $H=C/a$.

This proves that the first square charge has a unique regular solution,
$F=J_{02}$. It also identifies the complete polar ambiguity of that
equation. The spectral-projector square charge is
$J_{02}+\zeta^2/(2a)$. Replacing it by $J_{02}$ cancels its pole and
solves the first compatibility equation.

The linear functional
\[
 K((a))/K[[a]]\longrightarrow K,
 \qquad [h]\longmapsto\operatorname{Res}_{a=0}(h\,da)
\]
detects this ambiguity. The statement concerns the class of the function
in the quotient. It is not a residue of the symplectic form.

\subsection{Rigidity relative to five fixed traces}
\label{mht:frgl:sec:rigidity}

\begin{theorem}[Local rigidity]
\label{mht:frgl:thm:local-rigidity}
Every Lie map $\Phi:\mathfrak g_\tau\to K[[s,t,a,k]]$ agreeing with
$J$ on $\mathfrak b$ must satisfy $\Phi(q^mp^n)=J_{mn}$ for every
$m,n\geq0$. If $\zeta\ne0$, there is no such map, for any $\tau$.
\end{theorem}

\begin{proof}
The preceding calculation forces $\Phi(p^2)=J_{02}$.
Translation covariance gives
\[
 \Phi(q^3)=J_{30}+A(a,k),\qquad J_{30}=2s^3+6sa.
\]
Indeed the brackets with $q$ and $p$ say respectively that its $t$
derivative is zero and its $s$ derivative is $3J_{20}$.

The bracket with $q^2$ gives $A_k=0$. The bracket with $qp$ gives
$2aA'=3A$. For $A=\sum_{n\geq0}A_na^n$, the equation is
$(2n-3)A_n=0$ for every $n$. Every factor is nonzero in $K$, so
$A=0$. This proof permits arbitrary formal series and makes no degree
or finite-weight assumption.

The three quadratic Hamiltonians act by the usual linear changes of
the pair $(X,Y)$. In particular
\[
 \{J_{02}/2,J_{mn}\}=-mJ_{m-1,n+1}.
\]
The same formula holds in the source, since the source input $p^2/2$
has degree two. Hence knowing $\Phi(q^d)=J_{d0}$ determines every
degree-$d$ monomial by repeated bracketing with $p^2/2$.

The degree-three case is therefore fixed. Inductively, for $d\geq4$,
\[
 \mathcal B_\tau(q^3,q^{d-2}p)=3q^d,
 \qquad
 \{\operatorname{Tr}X^3,\operatorname{Tr}(X^{d-2}Y)\}
       =3\operatorname{Tr}X^d.
\]

The first formula has no higher Moyal term because its second input
has $p$ degree one. Its second input has total degree $d-1$, so the
induction applies. This fixes $q^d$, and then the whole degree-$d$
sector. Thus every component equals its symmetrized trace.

For the cubic pair one has
\[
 \{J_{30},J_{03}\}-9J_{22}=-3\zeta^2,
 \qquad \mathcal B_\tau(q^3,p^3)=9q^2p^2+\tfrac32\tau.
\]

Since $\Phi(1)=2$, compatibility forces $\tau=-\zeta^2$.
At that parameter the quartic residual is
\begin{equation}\label{mht:frgl:eq:local-residual}
 \{J_{40},J_{04}\}-16J_{33}-24\tau J_{11}
   =\frac{224}{5}\zeta^2 b
   =\frac{112}{5}\zeta^3+\frac{224}{5}\zeta^2ak.
\end{equation}
It is nonzero when $\zeta\ne0$.

For completeness, this identity can
be derived entirely in this chart by polarizing
\[
 \operatorname{Tr}(uX+vY)^n
   =2\sum_{j=0}^{\lfloor n/2\rfloor}
        \binom n{2j}(us+vt)^{n-2j}
          (a u^2+2buv+c v^2)^j
\]
and using \eqref{mht:frgl:eq:orbit}. It is also the rank-two specialization
of \eqref{mht:eq:framed-quartic-residual}. The
value at $s=t=a=k=0$ is $112\zeta^3/5$, proving the contradiction.
\end{proof}

The cyclic-word derivation of Proposition~\ref{mht:prop:framed-residual}
gives the same residual without choosing a chart.

The theorem excludes every correction that is regular as a formal germ
and fixes the stated five charges. It consequently excludes every
correction fixing the full degree-at-most-two scalar trace sector, or
the sector $K[q]\oplus K[q]p$. The correction need not be a trace, have bounded polynomial degree,
or be linear in framing invariants. Its order in $\zeta$ is unrestricted.

\subsection{Allowing the mixed quadratic to change}
\label{mht:frgl:sec:local-positive}

The preceding obstruction is relative to a specified sector. A regular
local map exists when the charge $qp$ is allowed to change. In fact all
such maps can be classified.

\begin{theorem}[Four-charge classification]
\label{mht:frgl:thm:four-charge}
Let $\Phi:\mathfrak g_\tau\to K[[s,t,a,k]]$ be a Lie map satisfying
\[
 \Phi(1)=2,\qquad\Phi(q)=2s,\qquad\Phi(p)=2t,
 \qquad\Phi(q^2)=2s^2+2a.
\]
Then $\tau=0$. There is a unique $h(a)\in K[[a]]$ such that
\begin{equation}\label{mht:frgl:eq:positive-map}
 \Phi(f)=f(s+\sqrt a,t+(k+h(a))\sqrt a)
          +f(s-\sqrt a,t-(k+h(a))\sqrt a).
\end{equation}
The expression on the right belongs to $K[[s,t,a,k]]$: its odd powers
of $\sqrt a$ cancel. Conversely, every $h\in K[[a]]$ defines such a
strict Poisson Lie map. If the images are regular functions on $U_+$,
then $h\in K[a]$.
\end{theorem}

\begin{proof}
Put $G=\Phi(qp)$. The brackets with $q,p,q^2$ give
\[
 G=2st+2ak+\alpha(a)
\]
for some $\alpha\in K[[a]]$. Set $\alpha_0=\alpha(0)$ and
$h=(\alpha-\alpha_0)/(2a)$. This is a formal power series.

The variable
$K_0=k+h(a)$ satisfies $\{a,K_0\}=1$. This identity is all that is
needed below. When $h(0)\ne0$, $K_0$ is not a parameter vanishing at
the original formal base point.

The equations for $F=\Phi(p^2)$ now give
\[
 F=2t^2+2ak^2+2\alpha(a)k+H(a),
 \qquad (aH)'=\alpha\alpha'.
\]
Regularity fixes the integration constant, so
\[
 H=\frac{\alpha(a)^2-\alpha_0^2}{2a},\qquad
 F=2t^2+2aK_0^2+2\alpha_0 K_0,
 \qquad G=2st+2aK_0+\alpha_0.
\]

In the variables $(s,t,a,K_0)$ these are precisely the scalar
quadratic traces for \eqref{mht:frgl:eq:local-gauge} with effective parameter
$\alpha_0$ in place of $\zeta$. The rigidity proof of
Theorem~\ref{mht:frgl:thm:local-rigidity} uses only their displayed brackets and
the power-series equation $2aA'=3A$, so it applies here as well. It forces all images
to be those effective symmetrized traces, and it forces
$\tau=-\alpha_0^2$. The quartic residual is
$224\alpha_0^2(\alpha_0/2+aK_0)/5$, so $\alpha_0=0$ and $\tau=0$.

When $\alpha_0=0$, the effective matrices commute. Their simultaneous
eigenvalues on $a\ne0$ are exactly the pairs in
\eqref{mht:frgl:eq:positive-map}.

For existence, on the double cover $a=r^2\ne0$, put
$x_\pm=s\pm r$ and $y_\pm=t\pm K_0r$. Direct differentiation gives
$\{x_i,y_j\}=\delta_{ij}$ and zero for the other pairs. The sum of
evaluation maps is therefore a Poisson Lie map.

Its values and bracket
defects are formal series containing only integral powers of $a$. The identity on the localization is an identity before localization,
because multiplication by $a$ is injective.

This proves existence and
strictness. The argument already proves uniqueness. In the polynomial
case $G$ is polynomial, hence $\alpha$ is polynomial and so is $h$.
\end{proof}

For $h=0$ the map has
\[
 \Phi(qp)=J_{11}-\zeta,\qquad
 \Phi(p^2)=J_{02}-2\zeta k.
\]
It is regular at the specified stable collision, retains the source
constant, and specializes to the ordinary trace map at $\zeta=0$.
It has no higher components because it is strict. The mixed quadratic charge has changed. Its quantum lift will be
constructed in Theorem~\ref{mht:frgl:thm:regular-local-quantum}, after the target
product has been defined.

\section{The second collision prevents global extension}
\label{mht:frgl:sec:second-collision}

\begin{theorem}[Global four-charge obstruction]
\label{mht:frgl:thm:global-four}
Suppose $\zeta\ne0$. No $K$-linear Lie map from $\mathfrak g_\tau$
to the global regular functions on the stable rank-two quotient can satisfy
$\Phi(1)=2$, $\Phi(q)=\operatorname{Tr}X$,
$\Phi(p)=\operatorname{Tr}Y$, and
$\Phi(q^2)=\operatorname{Tr}X^2$.
\end{theorem}

\begin{proof}
The restriction of a global map to $U_+$ is polynomial. By
Theorem~\ref{mht:frgl:thm:four-charge}, it must have $\tau=0$ and must be
\eqref{mht:frgl:eq:positive-map} for a polynomial $h(a)$.

Consider the independent regular family
\[
 X_-(a)=\begin{pmatrix}0&1\\a&0\end{pmatrix},\qquad
 Y_-(a)=\begin{pmatrix}0&0\\-\zeta&0\end{pmatrix},\qquad
 I_-=\binom10,\quad J_-=(2\zeta,0).
\]
It satisfies the same moment equation. It is stable for every $a$,
since $\det(I_-,Y_-I_-)=-\zeta$. When $a\ne0$ it also lies in
$U_+$. There its coordinates are $s=t=0$ and $k=-\zeta/a$, since
$b=\operatorname{Tr}(X_-Y_-)/2=-\zeta/2$.

Consequently the pullback of the square charge is
\[
 \Phi(p^2)|_{(X_-,Y_-,I_-,J_-)}
  =2a\left(-\frac\zeta a+h(a)\right)^2
  =\frac{2\zeta^2}{a}-4\zeta h(a)+2ah(a)^2.
\]
Its residue at $a=0$ is $2\zeta^2\ne0$. A global regular function
must pull back to a regular function on this affine line, a
contradiction. The residue is independent of every permitted $h$.
\end{proof}

The four fixed values are fewer than a full scalar quadratic sector.
The proof permits arbitrary changes of $qp$, $p^2$, all higher charges,
and all polynomial dependence on the framing variables after reduction.
It does not exclude maps changing $q^2$. For example the global map
\[
 f\longmapsto 2f\left(\frac{\operatorname{Tr}X}{2},
                       \frac{\operatorname{Tr}Y}{2}\right)
\]
is a strict Poisson Lie map retaining $1,q,p$. It changes $q^2$ and
does not specialize to the full trace map at $\zeta=0$. Thus neither
the center nor the first moments alone give the obstruction.

For any proposed Hamiltonian $L_\infty$ map into a dg Poisson model,
passing to degree-zero cohomology gives a strict Lie map. Suppose this cohomology is independently identified, as a Lie algebra,
with the stated regular functions. The obstruction then applies if
the four cohomology classes remain fixed. This is conditional on
that recipient comparison. No BRST-to-reduction identification is
being inferred solely from a shared name.

\subsection{A relative complex for the global polar obstruction}
For the remaining global constructions take $K=\mathbb C$.
The second family supplies a quotient-valued obstruction.
Let $L=\mathbb C[a,a^{-1}]$, $L_{\mathrm{reg}}=\mathbb C[a]$,
and use the two-term complex
\[
 \mathcal P^\bullet=(L_{\mathrm{reg}}\lhook\joinrel\longrightarrow L),
 \qquad |L_{\mathrm{reg}}|=0,\quad |L|=1.
\]
For every admitted polynomial $h$, the pullback of $\Phi_h(p^2)$
has the same class
\[
 [2\zeta^2/a]\in H^1(\mathcal P^\bullet)=L/L_{\mathrm{reg}}.
\]
It is closed because the complex has no degree-two term.

The well-defined functional
$[F]\mapsto\operatorname{Res}_{a=0}(F\,da)$ sends it to
$2\zeta^2$. Changing $h$ adds the regular polynomial
$-4\zeta h+2ah^2$ and hence does not alter the class. The formal
version uses $\mathbb C[[a]]\hookrightarrow\mathbb C((a))$.

This complex controls extension of the classified square charge
along the independent second family. It is a necessary relative
quotient of the global gluing problem, not an identification of its
full derived mapping space.
\subsection{The global quadratic correction problem}
\label{mht:frgl:sec:global-carrier}

Fix $\zeta\ne0$ and put $\kappa=\zeta/2$. The following presentation
is an isomorphism of algebraic Poisson varieties.

\subsubsection{An independent presentation of the quotient}

Write
\[
 s=\tfrac12\operatorname{Tr}X,\quad
 t=\tfrac12\operatorname{Tr}Y,\quad
 S=X-s\mathbf1_2,\quad T=Y-t\mathbf1_2,
\]
\[
 a=\tfrac12\operatorname{Tr}S^2,\quad
 b=\tfrac12\operatorname{Tr}(ST),\quad
 c=\tfrac12\operatorname{Tr}T^2.
\]
For traceless two-by-two matrices,
\[
 S^2=a\mathbf1_2,\quad T^2=c\mathbf1_2,\quad
 ST+TS=2b\mathbf1_2.
\]
Taking the trace of the framed equation gives $JI=2\zeta$.
Consequently $(\zeta\mathbf1_2-IJ)^2=\zeta^2\mathbf1_2$.
On the other hand $[S,T]^2=4(b^2-ac)\mathbf1_2$, whence
\[
 ac-b^2+\kappa^2=0.
\]

\begin{lemma}[The global rank-two carrier]\label{mht:frgl:lem:global-carrier}
The trace coordinates identify the quotient with
\[
 \mathbb A^2_{s,t}\times Q_\kappa,
 \qquad Q_\kappa=\operatorname{Spec}
 R_\kappa,\qquad
 R_\kappa=\mathbb C[a,b,c]/(ac-b^2+\kappa^2).
\]
Its Poisson bracket is
\[
 \{s,t\}=\tfrac12,\qquad
 \{a,b\}=a,\qquad\{a,c\}=2b,\qquad\{b,c\}=c,
\]
with all brackets between $s,t$ and $a,b,c$ equal to zero.
\end{lemma}

\begin{proof}
The columns $I,SI$ or $I,TI$ form a basis. Indeed, if both pairs
were dependent, $I$ would be a simultaneous eigenvector for $S,T$.
Then $[S,T]I=0$, whereas the framed equation gives
$[S,T]I=-\zeta I\ne0$.

On the first chart choose the basis so that $I=e_2$ and $SI=e_1$.
The relation and tracelessness then force
\[
 S=\begin{pmatrix}0&1\\a&0\end{pmatrix},\qquad
 T=\begin{pmatrix}0&k\\2\kappa+ak&0\end{pmatrix},\qquad
 J=(0,4\kappa).
\]
Thus
\[
 b=\kappa+ak,\qquad c=2\kappa k+ak^2.
\]
The chart is exactly $Q_\kappa\setminus D_-$, where
$D_-=(a=0,b=-\kappa)$. The inverse coordinate is
$k=(b-\kappa)/a=c/(b+\kappa)$. These formulas are used on their
respective open sets and agree on the overlap.

On the second chart take $I=e_2,TI=e_1$. Then
\[
 T=\begin{pmatrix}0&1\\c&0\end{pmatrix},\qquad
 S=\begin{pmatrix}0&\ell\\c\ell-2\kappa&0\end{pmatrix},
\]
\[
 b=c\ell-\kappa,\qquad a=c\ell^2-2\kappa\ell.
\]
This chart is the complement of $(c=0,b=\kappa)$.

The two indicated complements cover $Q_\kappa$, since
$\kappa\ne0$. Their matrices, together with the change of basis
on the overlap, establish the claimed isomorphism. The invariant
canonical matrix bracket gives the displayed relations directly.
\end{proof}

\subsubsection{Every regular quadratic correction is a shear}

Let a putative Lie map satisfy the four fixed charges. Brackets with
$2s=\Phi(q)$ and $2t=\Phi(p)$ imply that
\[
 \Phi(qp)=2st+H(a,b,c),\qquad
 \Phi(p^2)=2t^2+G(a,b,c).
\]
The brackets with $\Phi(q^2)/2=s^2+a$ give
\[
 \{a,H\}=2a,\qquad\{a,G\}=2H.
\]

\begin{lemma}[All global quadratic changes]\label{mht:frgl:lem:global-shears}
There is a unique $u(a)\in\mathbb C[a]$ such that
\[
 H=2b+a u(a),\qquad
 G=2c+2u(a)b+\tfrac a2u(a)^2.
\]
Conversely these formulas satisfy every scalar quadratic Lie
identity. They are the standard trace quadratics followed by the
Poisson automorphism
\[
 \alpha_u(a)=a,\qquad
 \alpha_u(b)=b+\tfrac a2u(a),\qquad
 \alpha_u(c)=c+u(a)b+\tfrac a4u(a)^2,
\]
which fixes $s,t$ and has inverse $\alpha_{-u}$.
\end{lemma}

\begin{proof}
Localizing at $a$ gives
$R_\kappa[a^{-1}]=\mathbb C[a,a^{-1},b]$ with
$c=(b^2-\kappa^2)/a$ and $\{a,-\}=a\partial_b$.
Its kernel is $\mathbb C[a,a^{-1}]$. Intersection with $R_\kappa$
is $\mathbb C[a]$. To see the latter intersection, restrict a Laurent
polynomial to the generic point of either divisor $D_\pm$
defined by $a=0,b=\pm\kappa$. A negative power of $a$ is a pole.

It follows that
\[
 H=2b+h(a),\qquad
 G=2c+\frac{2h(a)b}{a}+r(a),
\]
where $h\in\mathbb C[a]$ and $r\in\mathbb C[a,a^{-1}]$.

The expression $2h(a)b/a$ has at worst a simple pole. Hence
regularity of $G$ forces $r$ to have no pole of order greater
than one. The residues of $2h(a)b/a+r$ at $D_+$ and $D_-$ are
\[
 2\kappa h(0)+r_{-1},\qquad
 -2\kappa h(0)+r_{-1}.
\]

Both must vanish. Since $\kappa\ne0$, this gives $h(0)=0$ and
$r_{-1}=0$. Thus $h=au$ with $u\in\mathbb C[a]$ and $r$ is
polynomial.

Finally $\{H,G\}=2G$ is exactly
\[
 (ar)'=\left(\frac{a^2u^2}{2}\right)'.
\]
Polynomiality at $a=0$ eliminates the integration constant, giving
$r=au^2/2$.
Substitution verifies preservation of the quadric
equation and all three centered Poisson brackets under $\alpha_u$. Derivatives of $u(a)$ cancel in the bracket of its last two
images. Substitution also verifies $\alpha_{-u}\alpha_u=1$.
\end{proof}

The two residues are independent conditions. Removing either
divisor changes this classification: a constant $h(0)$ can then
be cancelled on the remaining branch. Thus the argument allows
arbitrary corrections and locates their global compatibility
obstruction, rather than merely measuring a pole of one selected
formula.

\subsubsection{Quadratic normalization forces all polynomial traces}

Compose a putative map with $\alpha_{-u}$. It now agrees with the
standard trace map on all polynomials of degree at most two.
Let $J_{mn}$ denote the symmetrized trace with $m$ copies of $X$
and $n$ copies of $Y$. Here $J_{00}=2$.

\begin{lemma}[Rigidity after quadratic normalization]\label{mht:frgl:lem:global-rigidity}
Suppose a Lie map to this global recipient agrees with the trace map
in degrees at most two. Then
$\Phi(q^mp^n)=J_{mn}$ for every $m,n\ge0$.
\end{lemma}

\begin{proof}
For a homogeneous cubic $f$, put $D(f)=\Phi(f)-J(f)$.
Brackets with $q,p$ show that $D(f)$ is annihilated by
$\partial_s,\partial_t$, hence belongs to $R_\kappa$.

Brackets with the quadratic polynomials show that $D$ intertwines
their $\mathfrak{sl}_2$ actions. In particular, the source cubic
weights under $\{qp,-\}$ are $-3,-1,1,3$. In $R_\kappa$,
the action $\{2b,-\}$ gives weights $-2,0,2$ on $a,b,c$.
Every monomial has an even weight, and the defining relation
has weight zero. Thus $R_\kappa$ has no odd weight eigenspace,
so $D=0$.

Proceed by induction on total degree $n\ge4$. The Moyal bracket
\[
 \mathcal B_\tau(q^3,q^{n-2}p)=3q^n
\]
has no higher correction, since the second input has $p$-degree
one. The inductive hypothesis and the canonical matrix bracket
give
\[
 3\Phi(q^n)
 =\{\operatorname{Tr}X^3,\operatorname{Tr}(X^{n-2}Y)\}
 =3\operatorname{Tr}X^n.
\]
Repeated bracketing with $p^2/2$ gives all homogeneous degree-$n$
monomials. On matrices the same operation polarizes
$\operatorname{Tr}(X-vY)^n$ and therefore gives the symmetrized
traces with exactly the same binomial factors. This completes
the induction.
\end{proof}

\section{An obstruction invariant under every formal first-component change}
\label{mht:frgl:sec:formal-all-corrections}

The formal assertion uses a specified reduction parameter. Put
\[
 A=\mathbb C[[\zeta]],\qquad
 \widehat R=\varprojlim_m
 \mathbb C[s,t,a,b,c,\zeta]/
  (ac-b^2+\zeta^2/4,\zeta^m),
\]
with
\[
 \{s,t\}=\tfrac12,\quad\{a,b\}=a,\quad
 \{a,c\}=2b,\quad\{b,c\}=c.
\]
The completion permits a polynomial coefficient at every order. The polynomial degrees need not be uniformly bounded.

The normal
form $b^2=ac+\zeta^2/4$ shows that this completion is $\zeta$-torsion
free. Its reduction is
\[
 R_0=\mathbb C[s,t]\otimes Q_0,
 \qquad Q_0=\mathbb C[a,b,c]/(ac-b^2).
\]
The algebra $\widehat R$ is specified independently of any candidate
Hamiltonian comparison.

Let $\tau\in\zeta A$ be the \emph{squared} Moyal parameter and let
\[
 \widehat{\mathfrak g}_\tau=
   (\mathbb C[q,p][[\zeta]],\mathcal B_\tau)
\]
have the coefficientwise bracket
\[
 \mathcal B_\tau(f,g)=\sum_{r\ge0}
 \frac{\tau^r}{2^{2r}(2r+1)!}
 m\left(\partial_q\otimes\partial_p-
        \partial_p\otimes\partial_q\right)^{2r+1}(f\otimes g).
\]
Thus $\tau=\theta^2$ in the usual divided Moyal commutator of
parameter $\theta$. The parameter $\tau$ deforms the source
bracket. The parameter $\zeta$ is the framed moment parameter. Choosing
$\tau(\zeta)$ is explicit additional data.

Write $J_\zeta(q^mp^n)=J_{mn}$ for the trace polarization formulas
on the rank-two carrier, extended continuously. At $\zeta=0$ the
cover
\[
 a=u^2,\quad b=uv,\quad c=v^2,\qquad\{u,v\}=\tfrac12
\]
gives
\[
 J_0(f)=f(s+u,t+v)+f(s-u,t-v).
\]
The two particle pairs have bracket one internally and zero with
each other, so $J_0$ is a Lie map for the Poisson source.

Consider a continuous $A$-linear Lie map
\[
 \Phi:\widehat{\mathfrak g}_\tau\longrightarrow\widehat R,
 \qquad \Phi(1)=2,\qquad\Phi\bmod\zeta=J_0.
\]
The affine and quadratic charges of any such map can be normalized
without altering its reduction.

\subsection{Formal normalization of the affine charges}

\begin{lemma}\label{mht:frgl:lem:affine-normalization}
A continuous Poisson automorphism of $\widehat R$, equal to the
identity modulo $\zeta$, conjugates any such candidate to one with
$\Phi(q)=2s$ and $\Phi(p)=2t$.
\end{lemma}

\begin{proof}
Suppose the affine charges have already been normalized modulo
$\zeta^n$, and write their order-$n$ errors as $u,v\in R_0$.
The central identity $\{\Phi(q),\Phi(p)\}=2$ gives
\[
 \partial_su+\partial_tv=0.
\]
Polynomial integration over the coefficient algebra $Q_0$ gives
$K\in R_0$ satisfying
\[
 \partial_tK=u,\qquad\partial_sK=-v.
\]
For example, integrate $u$ in $t$ and correct the resulting
$s$-derivative by a polynomial depending only on $s$.

Choose any lift $\widetilde K\in\widehat R$.
Conjugation by $\exp(\zeta^n\{\widetilde K,-\})$ changes the
two errors by $-K_t,K_s$, hence removes them. Its exponential
converges $\zeta$-adically, since its derivation raises the
$\zeta$-order by $n$. The successive automorphisms converge
coefficientwise. At every fixed coefficient only finitely many
steps contribute, so no degree bound on $K$ is required.
\end{proof}

\subsection{The quadratic normalization and its explicit contraction}

With the affine charges fixed, their brackets imply that
\[
 \Phi(q^2/2)-s^2,\qquad
 \Phi(qp)-2st,\qquad
 \Phi(p^2/2)-t^2
\]
belong to the centered algebra
$\widehat Q=\mathbb C[a,b,c][[\zeta]]/(ac-b^2+\zeta^2/4)$.
They give an $\mathfrak{sl}_2$ realization reducing to $a,2b,c$.

\begin{lemma}[The quadratic contraction]\label{mht:frgl:lem:quadratic-normalization}
For the standard Hamiltonian action,
\[
 H^1(\mathfrak{sl}_2,Q_0)=0.
\]
Every order-$n$ perturbation of the standard centered quadratic
realization is removed by a centered Hamiltonian conjugation
equal to the identity modulo $\zeta^n$.
\end{lemma}

\begin{proof}
Give $a,b,c$ degree one. The degree-$m$ component of $Q_0$
is the space of homogeneous polynomials of degree $2m$ in
$u,v$ under $a=u^2,b=uv,c=v^2$. It is the irreducible module
$V_{2m}$. Hence $Q_0=\bigoplus_{m\ge0}V_{2m}$. The sum is
algebraic. Every cocycle has image in a finite sum because it
has three polynomial values and the action preserves degree.

Use the conventional basis
$E=c,F=-a,H=2b$, satisfying
$[H,E]=2E,[H,F]=-2F,[E,F]=H$.
On a module write
\[
 C=\tfrac12H^2+EF+FE.
\]
It acts on $V_{2m}$ by $2m(m+1)$. For a cocycle
$\eta:\mathfrak{sl}_2\to V_{2m}$, put
\[
 S_\eta=\tfrac12H\cdot\eta(H)
          +E\cdot\eta(F)+F\cdot\eta(E).
\]

The invariance of
$\frac12H\otimes H+E\otimes F+F\otimes E$ and the cocycle
identity give
\[
 x\cdot S_\eta=C\eta(x)
 \qquad(x\in\mathfrak{sl}_2).
\]
Indeed, commuting $x$ past the first action produces two terms
that cancel by invariance, leaving the Casimir acting on
$\eta(x)$. Thus for $m>0$,
$\eta(x)=x\cdot(S_\eta/[2m(m+1)])$.
For $m=0$ the module is trivial, and a cocycle vanishes because
$\mathfrak{sl}_2=[\mathfrak{sl}_2,\mathfrak{sl}_2]$.
This proves the vanishing with an explicit polynomial primitive.

If a quadratic realization is standard modulo $\zeta^n$, its
order-$n$ difference is a cocycle for this module, by expansion
of its three Lie identities. Write it as $x\mapsto\{j(x),K\}$
using the preceding primitive. Conjugation by
$\exp(\zeta^n\{\widetilde K,-\})$, for a centered lift
$\widetilde K$, removes it and preserves $s,t$.
Successive conjugations converge as before.
\end{proof}

\begin{theorem}[Formal nonexistence with fixed reduction]
\label{mht:frgl:thm:formal-all-corrections}
There is no continuous $A$-linear Lie map
\[
 \Phi:\widehat{\mathfrak g}_\tau\longrightarrow\widehat R
\]
such that $\Phi(1)=2$ and $\Phi\bmod\zeta=J_0$.
Every higher first-component correction is admitted, including
corrections to $q,p,q^2,qp,p^2$.
The obstruction already occurs modulo $\zeta^3$.
\end{theorem}

\begin{proof}
After these two normalizations, every affine and scalar quadratic
charge is exactly its standard trace. The parity and pure-power
induction from Lemma~\ref{mht:frgl:lem:global-rigidity} work coefficientwise in
$\widehat Q$: $\{2b,-\}$ still has only even weights, and the
relation $ac-b^2+\zeta^2/4$ has weight zero. Therefore the entire
normalized candidate must equal $J_\zeta$.
The cubic identity forces $\tau=-\zeta^2$, and the quartic
defect is then
\[
 \frac{224}{5}\zeta^2b\ne0\quad\text{in }\widehat R.
\]
This proves Theorem~\ref{mht:frgl:thm:formal-all-corrections}.

For the finite-order obstruction, work modulo $\zeta^3$. Normalize
only at orders one and two, and apply rigidity through degree six.
The cubic relation then gives
$\tau\equiv-\zeta^2\pmod{\zeta^3}$, while the quartic residual
remains $(224/5)\zeta^2b\ne0\pmod{\zeta^3}$.
Thus the contradiction requires no convergence argument.
\end{proof}

\subsection{The controlling complex and the detector}

There is a finite-order cohomological form of the obstruction.
Let
\[
 \mathfrak g_0=(\mathbb C[q,p],\{-,-\}),\qquad
 \mathfrak h=\mathbb C[q,p]_{\le2},
\]
and let $R_0$ be a $\mathfrak g_0$-module by
$f\cdot r=\{J_0(f),r\}$.
The normalized correction complex is the relative Chevalley--
Eilenberg complex
\[
 C^\bullet(\mathfrak g_0,\mathfrak h;R_0).
\]
For a one-cochain,
\[
 (d\eta)(f,g)=\{J_0(f),\eta(g)\}
             -\{J_0(g),\eta(f)\}-\eta(\{f,g\}).
\]

The cubic identity after normalization modulo $\zeta^2$ first
forces $\tau\equiv0\pmod{\zeta^2}$: its trace defect is zero
to order one whereas its source scalar correction is $3\tau$.
It is therefore legitimate, for any remaining candidate, to set
\[
 \Omega_\zeta(f,g)
 =\{J_\zeta(f),J_\zeta(g)\}
      -J_\zeta(\mathcal B_\tau(f,g)),\qquad
 \omega_2=\zeta^{-2}\Omega_\zeta\bmod\zeta.
\]

The quotient is well-defined. The target relation starts in order two,
and the cubic test forces the same lower bound on the source correction.
Since $J_0$ is a Lie map, the curvature is divisible by $\zeta^2$.

The trace map is exactly $\mathfrak h$-equivariant, so $\omega_2$
is relative. The Jacobi identities in source and target give
$d\omega_2=0$ by taking the first nonzero coefficient of the
curvature identity.

The same argument below with zero curvature shows that every relative
one-cocycle is zero. In particular the normalized order-one correction
vanishes. The order-two equation is exactly $\omega_2+d\eta_2=0$,
with no quadratic contribution from an unretained order-one correction.

Suppose $\omega_2=-d\eta$ with $\eta$ relative. Its values on
cubics are centered, by translation equivariance, and have odd
$\{2b,-\}$ weights. Hence they vanish. For every $n\ge4$,
\[
 \omega_2(q^3,q^{n-2}p)=0.
\]
Induction in the displayed formula for $d\eta$, followed by
quadratic polarization, gives $\eta=0$. Thus a relative primitive
would force $\omega_2=0$.

If $\tau=\lambda\zeta^2+O(\zeta^3)$, the cubic pair gives
\[
 \omega_2(q^3,p^3)=-3(1+\lambda).
\]
If this scalar vanishes, then
\[
 \omega_2(q^4,p^4)=\frac{224}{5}b.
\]

Evaluation of this normalized representative at the point
$s=t=0,a=b=c=1$ of $\operatorname{Spec}R_0$ detects the latter
representative with value $224/5$. Its nontrivial cohomology class is established by the preceding exclusion of every relative primitive. Evaluation by itself is not a functional on arbitrary cohomology classes. This evaluates the residual Hamiltonian. It is not asserted to be a morphism of Poisson algebras.
The affine polynomial contraction and the explicit quadratic
Casimir contraction show why the nonzero relative class remains
an obstruction after all admitted formal first-component changes.

\subsection{Global functions on the completed stable framed family}
\label{mht:frgl:sec:geometric-formal-recipient}

The formal Hamiltonian recipient can be identified with the actual
global functions on the completed stable rank-two framed family.
This is an identification of Poisson algebras of functions. The
underlying stable family is not identified with its affinization
at the zero moment level.

Let $\mathscr X\to\mathbb A^1_\zeta$ be the stable quotient of
\[
 [X,Y]+IJ=\zeta\mathbf1_2
\]
in rank two. Write $s=\operatorname{Tr}X/2$ and
$t=\operatorname{Tr}Y/2$, and put $S=X-s\mathbf1_2$,
$T=Y-t\mathbf1_2$.
Stability makes at least one of $(SI,I)$ and $(TI,I)$ a basis.
If both pairs were dependent, $I$ would generate a common invariant line. This argument also applies
at $\zeta=0$.

The two basis choices give affine charts
\[
 U_+=\operatorname{Spec}\mathbb C[\zeta,s,t,a,k],\qquad
 U_-=\operatorname{Spec}\mathbb C[\zeta,s,t,c,\ell].
\]
In the first chart,
\[
 S=\begin{pmatrix}0&1\\a&0\end{pmatrix},\quad
 T=\begin{pmatrix}0&k\\\zeta+ak&0\end{pmatrix},\quad
 I=e_2,\quad J=(0,2\zeta).
\]
In the second,
\[
 T=\begin{pmatrix}0&1\\c&0\end{pmatrix},\quad
 S=\begin{pmatrix}0&\ell\\c\ell-\zeta&0\end{pmatrix},\quad
 I=e_2,\quad J=(0,2\zeta).
\]

On the overlap the change of basis is diagonal with entries
$k,1$. Hence
\begin{equation}\label{mht:frgl:eq:stable-family-gluing}
 \ell=k^{-1},\qquad c=ak^2+\zeta k,
 \qquad a=c\ell^2-\zeta\ell.
\end{equation}
In particular the overlap is $k\ne0$, equivalently $\ell\ne0$.

The invariant trace
\[
 b=\tfrac12\operatorname{Tr}(ST)
   =\zeta/2+ak=-\zeta/2+c\ell
\]
is globally regular and satisfies $ac-b^2+\zeta^2/4=0$.

For $m\ge1$ put
\[
 \mathscr X_m=\mathscr X\times_{\mathbb C[\zeta]}
                \mathbb C[\zeta]/(\zeta^m),
\qquad
 A_m=\mathbb C[s,t,\zeta]/(\zeta^m).
\]

\begin{theorem}[The geometric formal recipient]
\label{mht:frgl:thm:geometric-formal-recipient}
The trace functions induce a unital Poisson algebra isomorphism
\[
 A_m[a,b,c]/(ac-b^2+\zeta^2/4)
    \xrightarrow{\ \sim\ }
 \Gamma(\mathscr X_m,\mathcal O_{\mathscr X_m})
\]
for every $m\ge1$. Moreover
\[
 H^1(\mathscr X_m,\mathcal O_{\mathscr X_m})=0.
\]
For the $\zeta$-adic completion $\widehat{\mathscr X}$ along
$\mathscr X_1$, this gives the continuous unital Poisson
algebra isomorphism
\[
 \widehat R
   \xrightarrow{\ \sim\ }
 \Gamma(\widehat{\mathscr X},\mathcal O_{\widehat{\mathscr X}}),
\qquad
 \widehat R=\varprojlim_m
  \mathbb C[s,t,a,b,c,\zeta]/
   (ac-b^2+\zeta^2/4,\zeta^m).
\]
All these maps are compatible with reduction in $\zeta$.
\end{theorem}

\begin{proof}
The affine cover has affine overlap. Its quasicoherent cohomology
is therefore computed by the two-term \v Cech complex
\[
 C_m^0=A_m[k,a]\oplus A_m[\ell,c]
 \longrightarrow C_m^1=A_m[k,k^{-1},a],
\]
where the second summand is embedded by
$\ell=k^{-1}$, $c=k^2a+\zeta k$.
The differential is the difference of the two restrictions.
No field hypothesis on $A_m$ is used.

Filter each term by degree at most $n$ in the affine fiber
variable $a$, or $c$ on the second chart. The affine transition
\eqref{mht:frgl:eq:stable-family-gluing} preserves this filtration.
Call the resulting finite-degree complex $F_nC_m^\bullet$.
Its degree-$n$ quotient is
\[
 A_m[k]a^n\oplus A_m[k^{-1}]c^n
 \longrightarrow A_m[k,k^{-1}]a^n,
\]
with $c^n=k^{2n}a^n$ in the quotient. This is the \v Cech
complex of $\mathcal O_{\mathbb P^1_{A_m}}(2n)$, but its
cohomology follows directly from Laurent coefficients:
\[
 A_m[k]+k^{2n}A_m[k^{-1}]=A_m[k,k^{-1}],
\]
\[
 A_m[k]\cap k^{2n}A_m[k^{-1}]
       =\bigoplus_{j=0}^{2n}A_m k^j.
\]

The first equality gives $H^1(\operatorname{gr}_nC_m)=0$. The second gives a free basis $1,k,\ldots,k^{2n}$ for its
$H^0$ after factoring out $a^n$.

The filtrations are split on each chart and on the overlap,
so their successive short exact sequences are short exact
sequences of these two-term complexes. Starting with $n=0$,
the resulting elementary kernel--cokernel sequence gives
\[
 H^1(F_nC_m)=0,
\]
\[
 0\longrightarrow H^0(F_{n-1}C_m)
  \longrightarrow H^0(F_nC_m)
  \longrightarrow\bigoplus_{j=0}^{2n}A_m k^ja^n
  \longrightarrow0.
\]
Every polynomial cochain has finite fiber degree. Taking the
union over $n$, which preserves exact sequences of modules,
proves $H^1(C_m)=0$ and hence the asserted sheaf-cohomology
vanishing. It also determines the associated graded algebra
of $H^0(C_m)$.

Now filter
\[
 R_m=A_m[a,b,c]/(ac-b^2+\zeta^2/4)
\]
by total degree in $a,b,c$, with $A_m$ in degree zero.
The monic relation $b^2=ac+\zeta^2/4$ gives the unique
$A_m$-basis of monomials $a^ib^\epsilon c^j$,
$\epsilon\in\{0,1\}$. It follows directly that
\[
 \operatorname{gr}R_m=A_m[a,b,c]/(ac-b^2).
\]
The trace map $R_m\to H^0(C_m)$ respects the filtrations.
On the first chart its leading terms are $a,ak,ak^2$.
In degree $n$ it therefore sends the normal monomial basis
to the preceding Laurent-coefficient basis as follows:
\[
 a^{n-j}c^j\longmapsto a^n k^{2j}
       \quad(0\le j\le n),
\]
\[
 a^{n-1-j}bc^j\longmapsto a^n k^{2j+1}
       \quad(0\le j<n).
\]

Thus it is an isomorphism on every associated graded piece.

This proves surjectivity by induction: match a section in its leading
term, then match the remaining lower-degree section. A kernel
element would likewise have zero leading term and descend
to zero by the same induction. This proves the asserted
global-sections isomorphism over each nonreduced ring $A_m$.

The Poisson structure is the one induced by the canonical
framed reduction. On $U_+$ its brackets are
$\{s,t\}=1/2$, $\{a,k\}=1$, with $\zeta$ central.
Substitution of $b=\zeta/2+ak$ and $c=ak^2+\zeta k$
gives
\[
 \{a,b\}=a,\qquad\{a,c\}=2b,\qquad\{b,c\}=c,
\]
and all mixed brackets with $s,t$ vanish. The same identities
therefore hold globally and on every infinitesimal fiber.
The algebra isomorphism preserves them on its generators,
so it is a Poisson isomorphism.

A global section of the completed structure sheaf is a pair of
completed chart sections agreeing on the completed overlap. This
equalizer commutes with inverse limits. Equivalently, its
components are precisely compatible global sections on all
$\mathscr X_m$. Hence
\[
 \Gamma(\widehat{\mathscr X},\mathcal O_{\widehat{\mathscr X}})
 =\varprojlim_m\Gamma(\mathscr X_m,\mathcal O_{\mathscr X_m})
 =\varprojlim_m R_m.
\]

All brackets are continuous for this topology and the
finite-stage maps preserve them, proving the completed
Poisson assertion. This uses the displayed affine charts
and equalizers directly. No theorem on formal functions is
needed.
\end{proof}

The same calculation with coefficient ring
$\mathbb C[s,t,\zeta]$, without imposing $\zeta^m=0$,
also computes the ordinary global function ring of
$\mathscr X$. The special fiber $\mathscr X_1$ contains
the exceptional projective line $a=c=0$ in these charts. Its trace coordinates have $a=b=c=0$. Thus the theorem
does not identify the stable family or its probe categories
with the affine scheme of its global functions.

\begin{corollary}[The completed geometric target]\label{mht:frgl:cor:geometric-obstruction}
For every permitted $\tau$, no continuous Lie map
\[
 \widehat{\mathfrak g}_\tau\longrightarrow
 \Gamma(\widehat{\mathscr X},\mathcal O_{\widehat{\mathscr X}})
\]
sends $1$ to $2$ and reduces to $J_0$ modulo $\zeta$.
For this specific geometric recipient, no additional
comparison of global Poisson function algebras is missing.
\end{corollary}

\begin{proof}
Compose a putative map with the inverse of the continuous
Poisson isomorphism in
Theorem~\ref{mht:frgl:thm:geometric-formal-recipient}. It gives exactly
the map excluded by Theorem~\ref{mht:frgl:thm:formal-all-corrections}.
\end{proof}

\subsection{The meaning of the formal endpoint}

The classical reduction in this theorem is specifically reduction
modulo the framed parameter $\zeta$. It must not be substituted
silently for reduction in an independent matrix quantization
parameter $\hbar$. A quantum or dg comparison is excluded by this
theorem only when its relevant reduction induces the displayed
$\zeta$-adic map on zeroth cohomology. An isolated nonzero-level
comparison changing the four-charge sector, without this formal
extension to $\zeta=0$, remains outside the theorem.
\section{A regular quantum map on the collision chart}
\label{mht:frgl:sec:regular-quantum-chart}

The classical polynomial map on the chart with coordinates
$(s,t,a,k)$ has a regular quantum lift.  The construction uses
differential operators on even polynomials.  It specifies the
quantum recipient independently.

All elements in this section have cohomological degree zero and
the differential is zero.  Put $R=\mathbb C[\zeta,\hbar]$. The
moment parameter $\zeta$ is a spectator.  Define
\[
 D_R=R\langle s,a,t,k\rangle/
 \bigl([s,t]-\hbar/2,[a,k]-\hbar,
       [s,a],[s,k],[t,a],[t,k]\bigr).
\]

Its normal order is $s^i a^j t^u k^v$.  Thus $D_R$ is free over
$R$ by the ordered monomials and has the faithful differential
representation
\[
 s\mapsto s,\qquad a\mapsto a,\qquad
 t\mapsto-\frac{\hbar}{2}\partial_s,
 \qquad k\mapsto-\hbar\partial_a
 \quad\hbox{on }R[s,a].
\]
Faithfulness follows from the independence of differential
operators of distinct orders and the absence of $\hbar$-torsion.

Consequently $D_R/\hbar D_R$ is the independently specified
polynomial Poisson algebra with
$\{s,t\}=1/2$ and $\{a,k\}=1$.

For a variable $\tau$, write the polynomial Moyal Lie bracket as
\begin{equation}\label{mht:frgl:eq:quantum-squared-parameter}
 \mathcal B_\tau(f,g)=
 \sum_{j\geq0}\frac{\tau^j}{2^{2j}(2j+1)!}
 m\bigl(\partial_q\otimes\partial_p-
          \partial_p\otimes\partial_q\bigr)^{2j+1}
 (f\otimes g).
\end{equation}
The sum is finite on every pair of polynomials.  If the ordinary
Moyal product has parameter $\hbar$, its commutator divided by
$\hbar$ is $\mathcal B_{\hbar^2}$.  The constant polynomial $1$ is retained
as a central Lie element.

Introduce the two-particle Weyl algebra $W_R^{(2)}$ on
$(x_+,y_+),(x_-,y_-)$, with
$[x_\alpha,y_\beta]=\hbar\delta_{\alpha\beta}$ and the other
commutators zero.  Put
\[
 s=\frac{x_++x_-}{2},\quad \beta=\frac{x_+-x_-}{2},\quad
 t=\frac{y_++y_-}{2},\quad \varpi=y_+-y_-.
\]
Then $[s,t]=\hbar/2$, $[\beta,\varpi]=\hbar$, and the two pairs commute.
The transposition of particles fixes $s,t$ and sends
$(\beta,\varpi)$ to $(-\beta,-\varpi)$.  Its invariant algebra has ordered basis
$s^i t^j \beta^m \varpi^n$ with $m+n$ even.

\begin{lemma}[Polynomial radial restriction]\label{mht:frgl:lem:radial-restriction}
There is an $R$-linear unital algebra homomorphism
\[
 \rho:(W_R^{(2)})^{\mathfrak S_2}\longrightarrow D_R
\]
that fixes $s,t$ and is given on the other ordered factors by
\begin{equation}\label{mht:frgl:eq:radial-monomial}
 \rho(\beta^m \varpi^n)=
 \sum_{j=0}^{\lfloor n/2\rfloor}
 \frac{n!2^{n-2j}}{(n-2j)!j!}
 (-\hbar)^j a^{(m+n)/2-j} k^{n-j},
 \qquad m+n\text{ even}.
\end{equation}
Every exponent of $a$ and $\hbar$ on the right is nonnegative.
In particular,
\[
 \rho(\beta^2)=a,\qquad \rho(\beta\varpi)=2ak,\qquad
 \rho(\varpi^2)=4ak^2-2\hbar k.
\]
\end{lemma}

\begin{proof}
Represent $W_R^{(2)}$ on $R[s,\beta]$ by
$t=-(\hbar/2)\partial_s$ and $\varpi=-\hbar\partial_\beta$.
Its invariant operators preserve the even subspace
$R[s,\beta^2]$, identified with $R[s,a]$ by $a=\beta^2$.

Repeated differentiation gives, for every polynomial $F$,
\[
 \partial_\beta^n F(\beta^2)=
 \sum_{j=0}^{\lfloor n/2\rfloor}
 \frac{n!2^{n-2j}}{(n-2j)!j!}
 \beta^{n-2j}\partial_a^{n-j}F(a).
\]
This identity follows, for example, by extracting the coefficient
of $u^n$ in $F((\beta+u)^2)=F(\beta^2+2\beta u+u^2)$.

Multiplication by $\beta^m(-\hbar)^n$ and substitution
$k=-\hbar\partial_a$ give \eqref{mht:frgl:eq:radial-monomial}.
Since $m+n$ is even and $j\leq n/2$, one has
$(m+n)/2-j\geq0$.  Also $j\geq0$.

Restriction to a preserved subspace respects composition.
The formula proves that every restricted operator belongs to
$D_R$, so restriction is the asserted algebra homomorphism.
Faithfulness of the differential representation of $D_R$
turns equality of its operators into equality in $D_R$.
No inverse of $a$, $\zeta$ or $\hbar$ enters the construction.
\end{proof}

Let $\operatorname{Op}_\hbar$ denote Weyl symmetrization on one
particle.  The sum
\[
 \Delta(f)=\operatorname{Op}_\hbar f(x_+,y_+)
             +\operatorname{Op}_\hbar f(x_-,y_-)
\]
belongs to $(W_R^{(2)})^{\mathfrak S_2}$.  Set
$\mathcal J_\hbar(f)=\rho\Delta(f)$.

\begin{theorem}[Regular local quantum Hamiltonians]
\label{mht:frgl:thm:regular-local-quantum}
The map
\[
 \mathcal J_\hbar:
 (R[q,p],\mathcal B_{\hbar^2})\longrightarrow
 (D_R,\hbar^{-1}[-,-])
\]
is a strict Lie algebra homomorphism.  Division on the right
denotes the unique quotient of an already $\hbar$-divisible
commutator.  The map retains the constant central generator with
$\mathcal J_\hbar(1)=2$ and has the following ordered formulas:
\begin{align*}
 \mathcal J_\hbar(q)&=2s,&
 \mathcal J_\hbar(p)&=2t,\\
 \mathcal J_\hbar(q^2)&=2s^2+2a,&
 \mathcal J_\hbar(qp)&=2st+2ak-\hbar,\\
 \mathcal J_\hbar(p^2)&=2t^2+2ak^2-\hbar k.
\end{align*}

Its reduction modulo $\hbar$ is the polynomial Lie map
\begin{equation}\label{mht:frgl:eq:quantum-classical-radial}
 \mathcal J_0(f)=
 f(s+\sqrt a,t+k\sqrt a)+
 f(s-\sqrt a,t-k\sqrt a).
\end{equation}
The right side is interpreted as its even polynomial expansion
in $\sqrt a$.  In particular it is regular at $a=0$.
The construction extends continuously to the $\hbar$-adic and
$(\zeta,\hbar)$-adic completions, allowing polynomial support
to grow with parameter order.
\end{theorem}

\begin{proof}
Weyl symmetrization on one particle satisfies
\[
 [\operatorname{Op}_\hbar f,\operatorname{Op}_\hbar g]
 =\hbar\operatorname{Op}_\hbar \mathcal B_{\hbar^2}(f,g).
\]
This follows by expanding the exponential definition of the
Moyal product. The odd powers give
\eqref{mht:frgl:eq:quantum-squared-parameter}.  Operators from distinct
particles commute.  Adding their identities proves
$[\Delta(f),\Delta(g)]=\hbar\Delta(\mathcal B_{\hbar^2}(f,g))$.
Applying the algebra homomorphism $\rho$ proves the asserted
Lie identity.

The affine and quadratic formulas follow from
\[
 x_+^2+x_-^2=2s^2+2\beta^2,\qquad
 y_+^2+y_-^2=2t^2+\varpi^2/2,
\]
and
\[
 \frac12\sum_{\alpha=\pm}
       (x_\alpha y_\alpha+y_\alpha x_\alpha)
       =2st+\beta\varpi-\hbar.
\]
Reduction of \eqref{mht:frgl:eq:radial-monomial} modulo $\hbar$ keeps
only $j=0$, namely
$\beta^m \varpi^n\mapsto 2^n a^{(m+n)/2}k^n$.
This is precisely substitution $\beta=\sqrt a$, $\varpi=2k\sqrt a$
on invariant polynomials and proves
\eqref{mht:frgl:eq:quantum-classical-radial}.

For each polynomial input every displayed sum is finite, and
every term has nonnegative parameter order.  At fixed parameter
precision only finitely many input coefficients contribute.
Thus the maps are continuous on the stated completions, where
the algebraic identities hold coefficientwise.
\end{proof}

The theorem constructs Hamiltonian functions with their quantum
Lie bracket.  The map $\rho$ preserves associative products. The map
$\mathcal J_\hbar$ is a Lie map, since the product of two
particle sums also contains terms from different particles.
As a map into a differential graded Lie algebra concentrated in
degree zero, it has all higher Taylor components equal to zero.

The polynomial quantum recipient $D_R$ has been constructed
independently.  Section~\ref{mht:frgl:sec:native-local-quantum} constructs
its operation-preserving identification with a specified
normal-ordered quantum Hamiltonian reduction on the cyclic
matrix chart, including its moment character.  The classical
coordinates alone would not identify the quantum products. The differential-operator coordinate map supplies that step.

\subsection{Quantum reduction on the cyclic matrix chart}
\label{mht:frgl:sec:native-local-quantum}

The polynomial differential-operator recipient of
Section~\ref{mht:frgl:sec:regular-quantum-chart} is realized by a specified
local quantum Hamiltonian reduction.  The coordinate change is
regular at the stable collision and fixes the moment character.

Let
\[
 Q=\{(X,I)\in\operatorname{Mat}_2\times\mathbb C^2:
          \det((X-sI_2)I,I)\ne0\},\qquad s=\tfrac12\operatorname{Tr}X,
 \qquad a=s^2-\det X.
\]
Here $I$ is a column vector, and $I_2$ denotes the identity matrix.
Define $g=((X-sI_2)I,I)$ and
$C(s,a)=\left(\begin{smallmatrix}s&1\\a&s\end{smallmatrix}\right)$.
Cayley--Hamilton gives a $\operatorname{GL}_2$-equivariant isomorphism
\begin{equation}\label{mht:frgl:eq:cyclic-base-trivialization}
 Q\simeq \operatorname{GL}_2\times\mathbb A^2_{s,a},\qquad
 (g,s,a)\longmapsto (gC(s,a)g^{-1},ge_2).
\end{equation}
Indeed $(X-sI_2)^2=aI_2$ proves both the displayed formula for
$X$ and the formula for its inverse.  The group acts by left
multiplication on $g$ and fixes $s,a$.

Let $R=\mathbb C[\zeta,\hbar]$.  For a smooth affine
$\mathbb C$-variety $Z$, write
$\mathscr D_{\hbar,R}(Z)=
 R\otimes_{\mathbb C[\hbar]}\operatorname{Rees}_\hbar\mathscr D(Z)$.
In particular $\mathscr D_{\hbar,R}(Q)$ is generated by regular
functions and operators $\hbar v$ for algebraic vector fields
$v$, with $[\hbar v,f]=\hbar v(f)$.
In the original cotangent coordinates use
\[
 Y_{ij}=-\hbar\partial_{X_{ji}},\qquad
 J_i=-\hbar\partial_{I_i}.
\]
Thus $[X_{ij},Y_{kl}]=\hbar\delta_{il}\delta_{jk}$ and
$[I_i,J_j]=\hbar\delta_{ij}$.

For $\xi\in\mathfrak{gl}_2$, set
\[
 \rho_\xi(X)=[X,\xi],\qquad \rho_\xi(I)=-\xi I,
 \qquad \widehat\mu_\xi=\hbar\rho_\xi,
 \qquad \chi_\xi=\zeta\operatorname{Tr}\xi.
\]
This specifies the quantum ordering.  It has classical moment
$\operatorname{Tr}\xi([X,Y]+IJ)$ and obeys
\[
 [\widehat\mu_\xi,\widehat\mu_\eta]
   =\hbar\widehat\mu_{[\xi,\eta]},\qquad
 \chi_{[\xi,\eta]}=0.
\]

Equivalently, with matrix products taken in the displayed
Weyl algebra and the framing product ordered $I$ before $J$,
\[
 \widehat\mu_\xi=
 \operatorname{Tr}\xi\bigl(XY-YX-2\hbar I_2+IJ\bigr).
\]
The Rees operator is kept as an element of the algebra. It is
not replaced by an operator divided by $\hbar$.

Write $\mathcal I$ for the left ideal generated by
$\widehat\mu_\xi-\chi_\xi$, and define the local quantum
reduction by
\[
 \mathscr A_{\zeta,\hbar}(Q)
   =\bigl(\mathscr D_{\hbar,R}(Q)/\mathcal I\bigr)^{\operatorname{GL}_2}.
\]
The product is the product of invariant classes represented
in the right normalizer of the left ideal.

\begin{theorem}[An integral local quantum reduction]
\label{mht:frgl:thm:local-native-quantum-reduction}
Horizontal differential operators in
\eqref{mht:frgl:eq:cyclic-base-trivialization} induce a unital
$R$-algebra isomorphism
\begin{equation}\label{mht:frgl:eq:native-chart-quantum-isomorphism}
 D_R\xrightarrow{\ \sim\ }\mathscr A_{\zeta,\hbar}(Q),
 \qquad
 t=-\frac{\hbar}{2}\partial_s,\quad
 k=-\hbar\partial_a.
\end{equation}
It is defined without inverting $a$, $\hbar$ or $\zeta$ and
preserves the character $\chi_\xi=\zeta\operatorname{Tr}\xi$.

Its reduction modulo $\hbar$ is the canonical Poisson
coordinate map on the stable $X$-cyclic chart of the framed
quotient.  On the classical reduction, the gauge section is
\begin{equation}\label{mht:frgl:eq:companion-reduced-section}
 X=\begin{pmatrix}s&1\\a&s\end{pmatrix},\quad
 Y=\begin{pmatrix}t&k\\ak+\zeta&t\end{pmatrix},\quad
 I=e_2,\quad J=(0,2\zeta).
\end{equation}
\end{theorem}

\begin{proof}
The isomorphism of smooth affine bases
\eqref{mht:frgl:eq:cyclic-base-trivialization} induces an exact
isomorphism of Rees differential-operator algebras by the
chain rule.  Hence
\[
 \mathscr D_{\hbar,R}(Q)
  \simeq\mathscr D_{\hbar,R}(\operatorname{GL}_2)\otimes_R D_R.
\]

The moment operators act entirely on the first factor.
The vector fields $\rho_\xi(g)=-\xi g$ give a global frame
of its tangent bundle.  Ordered differential operators give
the PBW decomposition
\[
 \mathscr D_{\hbar,R}(\operatorname{GL}_2)
    \simeq \mathcal O_R(\operatorname{GL}_2)\otimes_R U_{\hbar,R}(\mathfrak{gl}_2)
\]
as an $R$-module, with coefficient functions on the left
and moment operators on the right.  Here
$\mathcal O_R(\operatorname{GL}_2)=R\otimes_{\mathbb C}\mathcal O(\operatorname{GL}_2)$ and
$U_{\hbar,R}(\mathfrak{gl}_2)$ is the $R$-algebra with
relations $uv-vu=\hbar[u,v]_{\mathfrak{gl}_2}$.

Sending $u_\xi$ to $\chi_\xi$ defines an algebra character
of $U_{\hbar,R}(\mathfrak{gl}_2)$, since the trace of a
commutator is zero.
Its kernel is generated by the $u_\xi-\chi_\xi$.
The PBW decomposition therefore identifies the quotient
by the left moment ideal, as a module, with
\[
 \mathscr D_{\hbar,R}(Q)/\mathcal I
       \simeq \mathcal O_R(\operatorname{GL}_2)\otimes_R D_R.
\]

This decomposition is equivariant for left translation:
the moment character is invariant.  A left-invariant
regular function on $\operatorname{GL}_2$ is constant, as follows by
evaluating the equality $f(hg)=f(g)$ at $g=1$.
Thus every invariant class has a unique horizontal
representative in $D_R$.

Horizontal operators commute with every moment operator.
They consequently normalize the left ideal and their
ordinary composition gives the reduction product.
The preceding uniqueness proves that this algebra map
is bijective.  All these arguments are over $R$.

For the classical statement, the canonical one-form is
$\lambda=\operatorname{Tr}(Y\,dX)+J\,dI$.
Put $\widetilde Y=g^{-1}Yg$,
$\widetilde J=Jg$ and $\vartheta=g^{-1}dg$.
Direct substitution into \eqref{mht:frgl:eq:cyclic-base-trivialization}
gives
\[
 \lambda=\operatorname{Tr}(\widetilde Y\,dC)
       +\operatorname{Tr}\bigl(
         [C,\widetilde Y]+e_2\widetilde J\bigr)\vartheta.
\]

On the moment level the second term is
$\zeta\operatorname{Tr}\vartheta$. Its exterior derivative
vanishes.  In the gauge $g=1$, the off-diagonal moment equations force equal
diagonal entries of $Y$. The diagonal equations give
$Y_{21}=ak+\zeta$ and $J=(0,2\zeta)$.
Thus \eqref{mht:frgl:eq:companion-reduced-section} holds and
$\lambda=2t\,ds+k\,da$ on the section.
The resulting brackets are
$\{s,t\}=1/2$, $\{a,k\}=1$, exactly the semiclassical
relations of $D_R$.
\end{proof}

The theorem fixes normal ordering through
$\widehat\mu_\xi=\hbar\rho_\xi$.
For Weyl ordering of the entire cotangent quadratic moment,
the divergence of $\rho_\xi$ on
$\operatorname{Mat}_2\times\mathbb C^2$ is
$-\operatorname{Tr}\xi$.  Consequently
\[
 \widehat\mu_\xi^{\mathrm{Weyl}}
       =\widehat\mu_\xi-\frac\hbar2\operatorname{Tr}\xi.
\]
Weyl-ordered reduction at a value $\zeta_{\mathrm W}$
therefore corresponds to the displayed normal convention
at $\zeta=\zeta_{\mathrm W}+\hbar/2$.
This is an explicit change of the quantum moment character,
not an equality of the two conventions with their numerical
parameters unchanged.

\begin{corollary}[A local quantum BFV arrow]
\label{mht:frgl:cor:local-quantum-bfv-arrow}
The horizontal inclusion of $\mathcal J_\hbar$ is a strict Hamiltonian
Lie map from $(R[q,p],\mathcal B_{\hbar^2})$ to the quantum BFV algebra
of the cotangent problem on $Q$.
It has central value $2$ and all higher Taylor components
zero.  Its image in degree-zero reduction is the map of
Theorem~\ref{mht:frgl:thm:regular-local-quantum} under
\eqref{mht:frgl:eq:native-chart-quantum-isomorphism}.
\end{corollary}

\begin{proof}
Choose a basis of $\mathfrak{gl}_2$ with structure constants
$f_{ab}^{\ c}$, and add ghosts $c^a$ of degree $1$ and
antighosts $\beta_a$ of degree $-1$, with
\[
 c^ac^b+c^bc^a=0,\qquad
 \beta_a\beta_b+\beta_b\beta_a=0,\qquad
 \beta_a c^b+c^b\beta_a=\hbar\delta_a^b.
\]
All ghosts and antighosts commute with the even differential
operator algebra.  The BFV carrier is the resulting graded
tensor product with $\mathscr D_{\hbar,R}(Q)$.

The charge and differential are
\[
 \Theta=c^a(\widehat\mu_a-\chi_a)
      -\frac12 f_{ab}^{\ c}c^ac^b\beta_c,
 \qquad Q_\hbar=\hbar^{-1}[\Theta,-]_{\mathrm s}.
\]

The moment representation identity and Jacobi imply
$\Theta^2=0$.  Concretely, in the faithful differential
representation with $\beta_a=\hbar\partial_{c^a}$,
the charge is $\hbar$ times the Chevalley differential
for $\rho$, minus multiplication by the degree-one
character $c^a\chi_a$.  That character is closed and
its square is zero.  Hence the charge squares to zero
as an operator and, by PBW faithfulness, in the algebra.

The normalized commutator is integral because the
algebra is supercommutative modulo $\hbar$.

Every horizontal element of $D_R$ commutes with the
moment operators and with the ghosts.  It is therefore
$Q_\hbar$-closed.  Horizontal inclusion is a strict
dg algebra map.  Composing with the strict Lie map
$\mathcal J_\hbar$ gives the assertion.
\end{proof}

No quasi-isomorphism from $D_R$ to the full absolute
BFV complex is asserted.  The corollary constructs
the actual closed Hamiltonians and their quantum
bracket.  It is local on the open cyclic chart $Q$. Extension to the entire stable framed quotient still
requires a global comparison.

\section{Descent of a first-component obstruction}
\label{mht:frgl:sec:bfv-obstruction-descent}

Let $\mathfrak g$ be a Lie algebra in degree zero and let
$(B,Q,\{-,-\})$ be a differential graded Poisson algebra whose
bracket has degree zero.  Suppose
$J_1:\mathfrak g\to Z^0(B)$ is linear.  Its binary defect is
\[
 D_{J_1}(f,g)=\{J_1(f),J_1(g)\}-J_1([f,g]).
\]
The differential is a derivation, so $QD_{J_1}=0$.

\begin{proposition}[The degree-zero detector]
\label{mht:frgl:prop:bfv-degree-zero-detector}
If $J_1$ is the first component of a Hamiltonian
$L_\infty$ morphism, then $[D_{J_1}(f,g)]=0$ in $H^0(B)$
for every pair $f,g$.  This necessary condition is unchanged
by exact modifications $J_1(f)\mapsto J_1(f)+QK(f)$ with
$|K(f)|=-1$.  A cochain map $e:B\to\mathbb C$ with
$e(D_{J_1}(f,g))\ne0$ therefore excludes such an extension.
\end{proposition}

\begin{proof}
Choose the Taylor convention
$QJ_2(f,g)=J_1([f,g])-\{J_1(f),J_1(g)\}$.
It makes the defect exact.  For $J'_1=J_1+QK$, the derivation
rule and $QJ_1=0$ give
\begin{align*}
 D_{J'_1}(f,g)-D_{J_1}(f,g)=Q\bigl(&
 \{K(f),J_1(g)\}+\{J_1(f),K(g)\}\\
 &+\{K(f),QK(g)\}-K([f,g])\bigr).
\end{align*}
Hence the cohomology class is unchanged.  A cochain map to a
complex concentrated in degree zero annihilates boundaries,
which proves the final assertion.
\end{proof}

For the framed matrix BFV algebra take even coordinates
$X,Y,I,J$ and ghost--antighost coordinates $(c,\Psi)$ of
degrees $(1,-1)$.  Put
$\mu=[X,Y]+IJ-\zeta I_N$ and use the canonical brackets
$\{X_{ij},Y_{kl}\}=\delta_{il}\delta_{jk}$,
$\{I_i,J_j\}=\delta_{ij}$ and
$\{c_{ij},\Psi_{kl}\}=\delta_{il}\delta_{jk}$.
The BFV differential has
\[
 QX=-[c,X],\quad QY=-[c,Y],\quad QI=-cI,
 \quad QJ=Jc,\quad Qc=-c^2,
 \quad Q\Psi=\mu-[c,\Psi].
\]

At every constraint point, evaluation of the even coordinates
there and of all ghosts and antighosts at zero is a cochain
map.  Indeed every displayed image evaluates to zero.
Evaluation is used only as a map of complexes. It is not a
Poisson map from a symplectic manifold to a point.

Nonexact modifications of $J_1$ are a different question.
The proposition alone gives no invariance under them.  Such
invariance requires a relative deformation calculation or a
normalization theorem covering the admitted modifications.

\begin{proposition}[Formal quantum specialization]
\label{mht:frgl:prop:quantum-specialization-obstruction}
Let $\widehat B_\hbar$ be a complete separated
$\mathbb C[[\hbar]]$-linear differential graded associative
algebra, with no $\hbar$-torsion, such that graded
commutators belong to $\hbar\widehat B_\hbar$.
Assume its normalized commutator and differential reduce to
the specified classical dg Poisson algebra $B_0$.
Let $\widehat{\mathfrak g}_\hbar$ be a specified continuous
$\mathbb C[[\hbar]]$-linear Lie source whose bracket reduces to
that of $\mathfrak g_0$. Every continuous
$\mathbb C[[\hbar]]$-linear Hamiltonian
$L_\infty$ map from $\widehat{\mathfrak g}_\hbar$ into
$(\widehat B_\hbar,Q_\hbar,\hbar^{-1}[-,-]_{\mathrm s})$
induces a Hamiltonian $L_\infty$ map into $B_0$ after
reduction modulo $\hbar$.  In particular, an obstruction to
every permitted classical first component excludes every
quantum lift with that classical reduction.
\end{proposition}

\begin{proof}
Each Taylor identity is an equality of continuous multilinear
maps over $\mathbb C[[\hbar]]$.  Reducing this equality
modulo $\hbar$ gives the corresponding classical Taylor
identity, with the same degrees and signs.  Applied to the
binary identity, this is the previous proposition.
No point evaluation on a nonzero-parameter Weyl algebra is
used.  Such an evaluation would not respect its defining
commutator relation.
\end{proof}

\begin{corollary}[Descent to a specified recipient]\label{mht:frgl:cor:recipient-obstruction}
Let $B$ be a specified dg Poisson recipient. Fix a Lie map
$\rho:H^0(B)\longrightarrow\mathcal O(\mathcal X_{2,\zeta})$
taking the scalar unit to the scalar unit.
At $\zeta\ne0$, if the composite of a proposed first component with
$\rho$ has the four values in Theorem~\ref{mht:frgl:thm:global-four}, no Hamiltonian
$L_\infty$ extension exists. Likewise a parameter-complete version
whose composite lands in $\widehat R$ and reduces to $J_0$ is
excluded by Theorem~\ref{mht:frgl:thm:formal-all-corrections}.
\end{corollary}

\begin{proof}
The binary Taylor equation makes the induced map to $H^0(B)$ a
strict Lie map. Composing with the specified Lie realization gives
the forbidden map. No higher Taylor component changes this necessary
equation. The same argument applies after reducing a quantum dg
algebra modulo $\hbar$ as in the preceding proposition. This takes
cohomology after reduction. It assumes no unproved interchange of
cohomology with specialization.
\end{proof}

In the polynomial framed BFV model there is a concrete degree-zero
realization to reduced functions. Restrict the ghost-free part of a
closed degree-zero function to the constraint locus. The coefficient
of a single ghost in its differential says that this restriction is
gauge invariant: the possible antighost contributions are multiples
of constraints. Boundaries restrict to zero for the same reason.

Moreover a degree-zero monomial involving ghosts contains at least
one ghost and one antighost. Bracketing two such terms leaves a ghost
factor, so its ghost-free restriction is zero. The induced bracket is
therefore the ordinary reduced bracket of the ghost-free terms.

This constructs the stated Lie realization without claiming a
quasi-isomorphism of the full BFV algebra, a formality theorem, or an
identification of its full center. Its coefficientwise completion is
admissible only with the stated regular reductions and continuous
restriction maps.
\section{The remaining global comparisons}
\label{mht:frgl:sec:remaining}

For the unframed matrix BRST complex, extend the formal action
across the discriminant. Retain stabilizer states and coherent
comparisons with the prescribed trace charges.
The framed problem must also account for the new local maps and
their global obstructions.

The first regular correction is not the final comparison. On the
$X$-cyclic chart, replacing the spectral square charge by $J_{02}$ solves
its first equation. Keeping the mixed quadratic then forces all traces
and gives the quartic obstruction.

Changing the mixed quadratic instead
produces the strict maps $\Phi_h$, and radial restriction gives an
integral quantum lift of $\Phi_0$. Theorem~\ref{mht:frgl:thm:local-native-quantum-reduction}
and Corollary~\ref{mht:frgl:cor:local-quantum-bfv-arrow} identify its local quantum
recipient and construct its strict BFV inclusion. The second collision detects every
$\Phi_h$ with the same polar coefficient $2\zeta^2$. Thus the missing
global step cannot be achieved by a different choice within that
classified family.

The framing-adic theorem has a wider correction class. Suppose the full trace reduction is fixed at $\zeta=0$. The
Hamiltonian contractions normalize every affine and quadratic correction.
The surviving relative class is nonzero, so changing those charges
cannot repair the comparison. Its failure at order two excludes an all-order formal lift.

Theorem~\ref{mht:frgl:thm:geometric-formal-recipient}
identifies the formal function algebra with the global functions of
the completed stable framed family at every coefficient stage and
after completion. Thus this particular geometric recipient comparison
is constructed. The theorem makes no assertion of numerical convergence
and requires none.

\begin{problem}[An isolated nonzero-level comparison]
\label{mht:frgl:qn:isolated}
Fix $\zeta\ne0$ and an independently specified global quantum reduction
$\mathcal B_{\zeta,\hbar}$, with its normalized commutator and its
classical reduction. Allow a global Hamiltonian comparison to change $q^2$ and every other
nonlinear charge, retaining $1\mapsto2$. Does one exist with the prescribed
classical reduction in $\hbar$, without a formal extension to $\zeta=0$?
Specify the admitted first-component changes and the operations of the
full source center that the comparison must preserve.
\end{problem}

Theorems~\ref{mht:frgl:thm:global-four} and
\ref{mht:frgl:thm:formal-all-corrections} settle respectively the four-charge
subproblem and the framing-adic subproblem. They do not identify
Problem~\ref{mht:frgl:qn:isolated} with either of them. The global
center-of-mass map already shows why the central constant and the two
linear moments alone are insufficient. A positive solution of that
problem must supply a new regular nonlinear moment sector. A negative
solution must give an obstruction invariant under those additional
changes. A rank-two obstruction excludes a construction claimed
uniformly for all ranks, but it does not alone prove an isolated
fixed-rank obstruction for every $N>2$.

\begin{problem}[Global quantum transitions]\label{mht:frgl:qn:quantum-transitions}
The local comparison with the specified normal-ordered Rees reduction
has been constructed in Theorem~\ref{mht:frgl:thm:local-native-quantum-reduction}.
For a different independently prepared quantum or derived recipient,
construct its operation-preserving comparison with this local model,
retaining the explicitly stated ordering and moment-level conversion.
Determine the transition maps between cyclic charts and their effect on
an admitted nonlinear Hamiltonian sector. The radial map cannot extend
globally with its four marked classical charges, by
Theorem~\ref{mht:frgl:thm:global-four}. A global comparison allowing different
charges must supply new compatible transition data. The horizontal
inclusion into the absolute BFV algebra is not asserted to be a
quasi-isomorphism.
\end{problem}

\begin{problem}[The full center]\label{mht:frgl:qn:full-center}
Fix an independently defined chiral or derived center containing the
polynomial Hamiltonian sector. Construct its restriction to the stated
Hamiltonian recipient, preserving all specified operations. Any such
restriction satisfying the hypotheses of the obstruction theorems is
excluded. If its polynomial restriction has different operations or a
different classical reduction, state and solve that new comparison.
Neither a scalar central generator nor a Hamiltonian Lie map identifies
an $E_2$ cochain object with a geometric center.
\end{problem}

Identifying any of these algebraic recipients with the mixed field
theory requires comparison maps for its fields, cyclic pairing,
support and boundary conditions. A Lie map of Hamiltonians alone
does not supply that physical comparison.
